% =============================================================
%  Sections added by the accelerometer-based prediction work.
%  Include from main.tex with % =============================================================
%  Sections added by the accelerometer-based prediction work.
%  Include from main.tex with % =============================================================
%  Sections added by the accelerometer-based prediction work.
%  Include from main.tex with % =============================================================
%  Sections added by the accelerometer-based prediction work.
%  Include from main.tex with \input{prediction_sections}.
%  Numerical values come from prediction/results_macros.tex which is
%  regenerated by scripts/build_prediction_report_assets.py.
% =============================================================

\section{Height-Difference Prediction}
\label{sec:height-prediction}

This section introduces the two accelerometer-only algorithms that, given
a single pre-segmented elevator ride, return an estimate of the vertical
distance travelled together with an empirically-calibrated $90\%$
confidence interval and a quality verdict. Both algorithms are
registered as first-class members of the \texttt{PredictAlgorithm} enum
under \texttt{src/prediction/algorithms/} and instantiated through the
\texttt{Predictor} dispatcher (\S\ref{sec:predictor-api}), so that the
downstream inference code is algorithm-agnostic.

\subsection{Predictor API}
\label{sec:predictor-api}

Every algorithm implements three public methods (the signatures live in
\texttt{\small src/prediction/\allowbreak{}algorithms/\allowbreak{}predictor.py}):

\begin{itemize}
  \item \texttt{predict(data, phone\_model, pre, post)} $\;\to\;$
        \texttt{PredictionOutput} with fields
        \{\texttt{height\_diff, ci\_half\_width, theoretical\_sigma,
        accepted, quality\_score, reject\_reason, meta}\}. The
        \texttt{PredictionOutput} dataclass is defined in
        \texttt{\small src/prediction/\allowbreak{}algorithms/\allowbreak{}common/\allowbreak{}types.py}.
  \item \texttt{forward(data)} $\;\to\;$ \texttt{float}, the legacy
        float-in-float-out path used by the barometer baseline and
        external code.
  \item \texttt{calibrate(samples)} — fits the conformal multiplier on
        a set of labelled train samples and persists the checkpoint
        next to the calibration JSON. No-op for the barometer
        baseline.
\end{itemize}

Both accelerometer algorithms share a lightweight \texttt{common/}
package that provides: a phone-model $\to$ sensor-chip $\to$ noise
lookup (\texttt{common/noise\_db.py}), a split-conformal calibrator
(\texttt{common/conformal.py}), and a handful of vector helpers
(gravity estimation, vertical projection, ZUPT integration). Everything
else — quality filter, theoretical-$\sigma$ model, fitter — is
per-algorithm.

\subsection{Algorithm A: ZUPT Double-Integration}
\label{sec:algo-zupt}

Zero-Velocity Update (ZUPT) double-integration exploits the two facts
highlighted in \S\ref{sec:elevator-kinematics}: the cabin starts and
ends at rest, so in the absence of sensor noise the velocity
$v(t)=\int_0^t a(\tau)\,d\tau$ must satisfy $v(0)=v(T)=0$. We enforce
this boundary condition explicitly and apply a linear drift correction
to absorb the slow-varying bias that the static-bias model of the
accelerometer cannot cancel:
\begin{equation}
  v(t) = \int_0^t a(\tau)\,d\tau \;-\; \frac{t}{T}\,\int_0^T a(\tau)\,d\tau,
  \qquad
  \Delta h = \int_0^T v(t)\,dt.
  \label{eq:zupt-corr}
\end{equation}
The algorithm operates on the vertical component of acceleration,
$a_\mathrm{vert}(t)$, obtained by one of three strategies, chosen
automatically per segment based on pre/post-ride gravity stability
(the reasoning behind the cascade is in \texttt{zupt\_accel/estimator.py}):
\begin{enumerate}
  \item[(i)]   \emph{Pre+post projection.} When both pre- and post-ride
               stationary windows have $|{\bf g}|\in[8,12]$~m/s$^2$ and
               window-to-window standard deviation below a threshold,
               we compute a weighted mean gravity vector
               $\hat{\bf g}=(w_1{\bf g}_{\rm pre}+w_2{\bf g}_{\rm post})/(w_1+w_2)$
               with weights $w_i=1/\sigma_i$, and project
               $a_\mathrm{vert}=\mathbf{a}\cdot\hat{\bf g}/|\hat{\bf g}|-|\hat{\bf g}|$.
  \item[(ii)]  \emph{One-sided projection.} If only one of the windows
               is stable, its gravity vector is used.
  \item[(iii)] \emph{Magnitude fallback.}
               $a_\mathrm{vert}=|{\bf a}|-\bar{|{\bf a}|}$. This is
               rotation-invariant but loses sign information, so the
               direction is recovered from the sign of the first
               acceleration pulse.
\end{enumerate}
The active-motion window $[t_s,t_e]$ is located by thresholding the
moving-average of $|a_\mathrm{vert}|$ at $0.1$~m/s$^2$; acceleration
outside the window is zeroed before integration. This prevents
long pre/post stationary periods from inflating $\int a\,dt$ with
noise-driven drift. The number of samples inside the window,
$N=t_e-t_s+1$, is the key scale for the theoretical error analysis
below (\S\ref{sec:theoretical-sigma-zupt}).

\subsection{Algorithm B: 7-Step S-Curve Fitting}
\label{sec:algo-scurve}

Rather than integrating blindly, Algorithm B fits the entire ride to
the 7-phase S-curve template of \S\ref{sec:elevator-kinematics}. The
fitted parameters are $\theta=(j_{\max}, a_{\max}, v_{\max}, H, t_0)$
where $t_0$ is the motion-onset offset within the segment; direction
is recovered separately from ZUPT drift sign. The cost function is
\textbf{the velocity-domain residual}:
\begin{equation}
  \mathcal{L}(\theta) =
  \sum_{i=1}^{N}\left(\frac{v_i - v_\theta(t_i)}{\sigma_v}\right)^2
  \;+\;
  \lambda\sum_{k}\left(\frac{\theta_k-\mu_k}{\sigma_k}\right)^2,
  \label{eq:scurve-loss}
\end{equation}
where $v_i$ is the measured velocity (ZUPT-integrated, drift-corrected,
low-pass-filtered at $3$~Hz), $v_\theta(t)$ is the template velocity
derived analytically from $\theta$, $\sigma_v$ is a robust MAD estimate
of the velocity noise, and the second term is a Bayesian prior
regulariser whose parameters $(\mu_k,\sigma_k)$ are taken from ISO
18738 / CIBSE Guide D building-class priors. Moving to the velocity
domain yields a roughly 10$\times$ SNR improvement over the raw
acceleration domain because integration acts as a first-order
low-pass, and the velocity template (one smooth bump) is much more
distinctive than the sharp-edged acceleration template.

Fitting uses trust-region reflective NLS (\texttt{scipy.optimize.least\_squares},
method \texttt{trf}), warm-started from a 2-phase grid search in
distance and motion-onset that uses the prior-mean kinematic values.
Per-segment uncertainty on the distance parameter comes from the
Cram\'er--Rao bound
\begin{equation}
  \Sigma \;=\; \left(\mathbf{J}^\top \mathbf{J}/\sigma_v^2 + \Sigma_\mathrm{prior}^{-1}\right)^{-1},
  \qquad
  \sigma_{\rm CRB} = \sqrt{\Sigma_{HH}},
  \label{eq:scurve-fim}
\end{equation}
where $\mathbf{J}_{i,k}=\partial v_\theta(t_i)/\partial\theta_k$ and
$\Sigma_\mathrm{prior}^{-1}$ is the diagonal prior precision. The CRB
is optimistic because the ideal-template assumption leaks real ride
irregularity into the residuals, so we compose three variance terms
into the per-segment $\sigma$:
\[
  \sigma^2 \;=\; \bigl(f_\mathrm{safe}\,\sigma_{\rm CRB}\bigr)^2
  \;+\; \bigl(k_\mathrm{rel}\,|\hat{H}|\bigr)^2
  \;+\; \bigl(\tfrac{1}{2}|\hat{H}-\hat{H}_\mathrm{ZUPT}|\bigr)^2.
\]
The first term is the FIM result with a safety factor
$f_\mathrm{safe}$; the second is a per-meter floor encoding the
empirical observation that longer rides drift proportionally more; the
third is a cross-check penalty proportional to the disagreement
between the S-curve $\hat{H}$ and the independent ZUPT integrator's
$\hat{H}_\mathrm{ZUPT}$.

\subsection{Quality Filter}
\label{sec:quality-filter}

Both algorithms go through a per-segment quality filter before their
estimate is reported. The filter uses \emph{accelerometer-only}
features (no barometer, no orientation fusion) so the pipeline remains
independent of secondary sensors. Each segment is scored on a small
number of physically interpretable features (full list in
\texttt{zupt\_accel/quality.py} and the S-curve estimator):
\begin{itemize}
  \item \textbf{Pre/post gravity magnitude} and \textbf{window
        stability} — the phone should rest between $[8,12]$~m/s$^2$
        with intra-window $\sigma_{|{\bf g}|}\lesssim 1$~m/s$^2$.
  \item \textbf{Pre-vs-post gravity angle} — large rotation of the
        phone body frame breaks the gravity-projection step.
  \item \textbf{During-ride gravity drift} — max angle swept across
        1-s chunks; bounded by $15^\circ$ before we reject.
  \item \textbf{Max transient peak} — impacts above $8$~m/s$^2$
        dominate the integral.
  \item \textbf{Active-motion fraction} — fraction of samples with
        $|a_\mathrm{vert}|>0.1$~m/s$^2$; very low fractions on long
        segments indicate a non-ride.
  \item \textbf{ZUPT-vs-S-curve direction agreement} (S-curve only).
\end{itemize}
The features are combined into a continuous score in $[0,10]$ via
weights tuned on the training set; on top of the score, each feature
also has a hard threshold that triggers an outright rejection with a
human-readable reason (e.g.\ \texttt{"ride\_drift\_22deg"}). The
quality filter is a \emph{filter}, not a predictor: it is designed to
let through almost all the clean data even at the cost of letting
through a handful of dirty segments, because the conformal CI is what
provides the coverage guarantee; the filter's job is to catch the
pathological cases where the per-segment $\sigma$ model is guaranteed
to under-predict the real error.

\subsection{Theoretical Per-Segment $\sigma$}
\label{sec:theoretical-sigma-zupt}

The classical ZUPT analysis treats acceleration noise as white with
standard deviation $\sigma_a$ (datasheet value, from the sensor chip
— see \texttt{common/noise\_db.py}). Under a linear drift correction,
the variance of the end-point position estimate is
\begin{equation}
  \mathrm{Var}(\hat{H}) \;=\; \sigma_a^2\,\Delta t^{4}\,\frac{N^3}{12},
  \qquad
  \sigma_\mathrm{white} \;=\; \sigma_a\,\Delta t^{2}\,\sqrt{\frac{N^3}{12}},
  \label{eq:zupt-var}
\end{equation}
where $N$ is the number of active-motion samples and $\Delta t$ is
the sampling interval (50~Hz phones $\Rightarrow \Delta t = 20$~ms).
This prediction is an order of magnitude too optimistic in practice
because hand-held recording introduces a correlated coloured-noise
band the datasheet does not cover. We capture the coloured-noise
contribution with a second term whose scale grows linearly with the
quality-filter gravity-drift angle (up to $3\times$), and a third term
proportional to the predicted distance itself,
\begin{equation}
  \sigma^2_{\rm ZUPT} \;=\;
  \sigma_\mathrm{white}^2
  + \sigma_\mathrm{mech}^2(\rho)
  + k_\mathrm{rel}^2\,|\hat{H}|^2,
  \label{eq:zupt-sigma-total}
\end{equation}
with $\rho$ the drift angle in degrees. The per-meter relative floor
$k_\mathrm{rel}$ absorbs the long-ride error growth the physics model
misses. All three components are cheap to evaluate per segment, so
the theoretical $\sigma$ varies from ride to ride even before the
conformal layer enters.

The S-curve estimator uses the analogous composition around the CRB
of \eqref{eq:scurve-fim}, as described in \S\ref{sec:algo-scurve}.

\subsection{Conformal Calibration}
\label{sec:conformal-calibration}

Neither the white-noise model \eqref{eq:zupt-var} nor the CRB
\eqref{eq:scurve-fim} is guaranteed to be calibrated on real data.
We therefore close the loop with a \emph{split-conformal} step on
the training set: for every calibration sample $i$ compute the
non-conformity score
\begin{equation}
  r_i \;=\; \frac{|\hat{H}_i - H_i|}{\sigma_i},
  \label{eq:nonconformity}
\end{equation}
and take the $(1-\alpha)(n+1)/n$ empirical quantile of the
$\{r_i\}$ as the conformal multiplier $k$. The reported CI is then
$[\hat{H}\pm k\sigma]$, subject to absolute floor and cap guards
against pathological $\sigma$. By the standard split-conformal
result this guarantees
\begin{equation}
  \mathbb{P}\bigl(\,|\hat{H}-H|\le k\sigma\,\bigr) \;\ge\; 1-\alpha
  \label{eq:conformal-coverage}
\end{equation}
for \emph{exchangeable} data. Exchangeability is broken across the
clean-vs-dirty split so we calibrate on clean \emph{and}
quality-filter-accepted samples only, which is our proxy for the
exchangeable subset. The calibration JSON checkpoints live under
\texttt{\small src/data/\allowbreak{}structuredData/\allowbreak{}test\_results/\allowbreak{}prediction/\allowbreak{}train/}.

\section{Evaluation Protocol}
\label{sec:evaluation-protocol}

\paragraph{Train / test split.} The train/test split is
\emph{location-based}, not random: every experiment recorded at
\texttt{beitYitzchakiRaanana} is held out for final test, and the
remaining experiments (Millenium, Acro, Beit Mansour, Bar-Ilan~2,
Haari) form the train set. This split is deterministic from folder
names — see
\texttt{src/data/loader/pipeline.classify\_experiment\_type}. Because
the held-out location was never visited during training, a test-set
win is a statement about \textbf{cross-building generalisation},
which is the property we actually care about.

\paragraph{Clean vs.\ dirty.} Every elevator segment in
\texttt{gt.csv} carries a \texttt{signalClearRecording} Boolean which
is \texttt{True} when the experimenter marked the ride as a
clean recording during GT editing and \texttt{False} when the signal
is known to be corrupted (e.g.\ the experimenter stopping and
starting the logger mid-ride). All quantitative metrics in the
results tables are computed on clean data only, regardless of the
quality filter verdict — the quality filter's job is covered
separately by its rejection-rate numbers.

\paragraph{Metrics.}
For each algorithm we report, on the clean subset:
\begin{itemize}
  \item \textbf{Coverage} $=\mathbb{P}(|\hat{H}-H|\le w)$ where $w$ is
        the reported CI half-width; target $\ge 90\%$.
  \item \textbf{Median $|e|$} and \textbf{MAE}: $|e|=|\hat{H}-H|$.
  \item \textbf{Median CI half-width $\tilde{w}$}: a tighter
        $\tilde{w}$ is better as long as coverage stays above target.
  \item \textbf{Fraction within one floor} ($\le 1.5$~m):
        an absolute usability number.
  \item \textbf{Dirty-rejection rate}: fraction of
        \texttt{signalClearRecording=False} segments rejected by the
        quality filter.
\end{itemize}

\paragraph{Single-shot test pass.} Until the final run, only train
metrics are inspected and iterated upon. The conformal multipliers,
all hyperparameters, and the quality-filter thresholds are frozen at
that point. The blind-test pass is then run exactly once via
\texttt{\small python scripts/\allowbreak{}run\_prediction\_evaluation.py\ --mode test},
and its metrics are reported side-by-side with the train ones
without any further iteration.

\section{Results}
\label{sec:prediction-results}

Figures and tables in this section are produced by
\texttt{scripts/run\_prediction\_evaluation.py} and
\texttt{scripts/build\_prediction\_report\_assets.py}.

\subsection{Train-set headline numbers}

Table~\ref{tab:headline-train} summarises the train-set metrics. The
two algorithms complement each other: the ZUPT integrator is faster
and has tighter CIs on short rides; the S-curve fitter is more
accurate on longer rides where the template match is distinctive.

\begin{table}[H]
\centering
\input{figures/prediction/results_table_train}
\caption{Headline metrics on the \textbf{train} split (clean
segments). ``Cov'' columns are $90\%$-CI coverage; $\tilde{|e|}$
and $\tilde{w}$ are the medians.}
\label{tab:headline-train}
\end{table}

\subsection{Blind test results}
\label{sec:blind-test}

Table~\ref{tab:headline-test} reports the final blind-test metrics.
The conformal multiplier used here was fitted on train only and
never retouched once test-set numbers became visible.

\begin{table}[H]
\centering
\input{figures/prediction/results_table_test}
\caption{Headline metrics on the \textbf{test} split
(\texttt{beitYitzchaki}). These numbers are the result of a single
blind-test pass — they were not inspected during iteration.}
\label{tab:headline-test}
\end{table}

\subsection{Figures}

\begin{figure}[H]
  \centering
  \begin{subfigure}{0.48\linewidth}
    \includegraphics[width=\linewidth]{prediction/train_zupt_scatter.png}
    \caption{ZUPT, train}
  \end{subfigure}\hfill
  \begin{subfigure}{0.48\linewidth}
    \includegraphics[width=\linewidth]{prediction/train_scurve_scatter.png}
    \caption{S-curve, train}
  \end{subfigure}
  \caption{Predicted vs.\ true $\Delta h$ on the train split. Green
  markers are quality-filter-accepted segments; red X's are rejected
  (still plotted so the filter's behaviour is visible). The dotted
  band is the $\pm 1.5$~m (one floor) target.}
  \label{fig:train-scatter}
\end{figure}

\begin{figure}[H]
  \centering
  \begin{subfigure}{0.48\linewidth}
    \includegraphics[width=\linewidth]{prediction/train_zupt_cdf.png}
    \caption{ZUPT, train}
  \end{subfigure}\hfill
  \begin{subfigure}{0.48\linewidth}
    \includegraphics[width=\linewidth]{prediction/train_scurve_cdf.png}
    \caption{S-curve, train}
  \end{subfigure}
  \caption{Empirical CDF of $|e|$. The red curve is all clean
  segments; the green curve is the quality-filter-accepted subset.
  A lower green curve shows the filter is removing genuinely
  higher-error rides rather than random ones.}
  \label{fig:train-cdf}
\end{figure}

\begin{figure}[H]
  \centering
  \begin{subfigure}{0.48\linewidth}
    \includegraphics[width=\linewidth]{prediction/train_zupt_ci.png}
    \caption{ZUPT, train}
  \end{subfigure}\hfill
  \begin{subfigure}{0.48\linewidth}
    \includegraphics[width=\linewidth]{prediction/train_scurve_ci.png}
    \caption{S-curve, train}
  \end{subfigure}
  \caption{Per-segment CI plots. Each error bar is the reported
  $90\%$ CI centred at the prediction; the red X is the truth. Green
  background = segment is covered, red = miss. Segments are sorted
  by true $\Delta h$ so the per-segment CI adaptivity can be read off
  directly.}
  \label{fig:train-ci}
\end{figure}

\begin{figure}[H]
  \centering
  \begin{subfigure}{0.48\linewidth}
    \includegraphics[width=\linewidth]{prediction/train_zupt_cov_bins.png}
    \caption{ZUPT, train}
  \end{subfigure}\hfill
  \begin{subfigure}{0.48\linewidth}
    \includegraphics[width=\linewidth]{prediction/train_scurve_cov_bins.png}
    \caption{S-curve, train}
  \end{subfigure}
  \caption{Coverage and MAE broken down by $|\Delta h|$ bin. A healthy
  calibrator keeps the green bars close to the dotted $90\%$ target
  across bins — uniform under-coverage in the long-distance bin would
  mean the relative $\sigma$ term is too small.}
  \label{fig:train-cov-bins}
\end{figure}

\begin{figure}[H]
  \centering
  \includegraphics[width=0.75\linewidth]{prediction/train_fig_compare_algorithms_train.png}
  \caption{Direct comparison of the two algorithms' error CDFs on the
  train split (clean only). The closer the curve pushes toward the
  top-left, the better.}
  \label{fig:train-compare}
\end{figure}

\begin{figure}[H]
  \centering
  \begin{subfigure}{0.48\linewidth}
    \includegraphics[width=\linewidth]{prediction/train_zupt_reliability.png}
    \caption{ZUPT}
  \end{subfigure}\hfill
  \begin{subfigure}{0.48\linewidth}
    \includegraphics[width=\linewidth]{prediction/train_scurve_reliability.png}
    \caption{S-curve}
  \end{subfigure}
  \caption{Reliability diagrams on the train split. Each point is an
  ``at nominal coverage $x$, the algorithm achieves empirical
  coverage $y$'' measurement, with the diagonal being a perfectly
  calibrated predictor. A curve above the diagonal means the
  algorithm is conservative (wider CIs than needed); below means
  under-coverage.}
  \label{fig:train-reliability}
\end{figure}

\begin{figure}[H]
  \centering
  \begin{subfigure}{0.48\linewidth}
    \includegraphics[width=\linewidth]{prediction/train_zupt_signed_err.png}
    \caption{ZUPT}
  \end{subfigure}\hfill
  \begin{subfigure}{0.48\linewidth}
    \includegraphics[width=\linewidth]{prediction/train_scurve_signed_err.png}
    \caption{S-curve}
  \end{subfigure}
  \caption{Signed error $\hat{\Delta h}-\Delta h$ vs.\ true $\Delta h$
  on train. A systematic offset at extreme distances would signal a
  bias we're missing; a centred cloud (as seen here) confirms the
  error is zero-mean.}
  \label{fig:train-signed}
\end{figure}

\begin{figure}[H]
  \centering
  \begin{subfigure}{0.48\linewidth}
    \includegraphics[width=\linewidth]{prediction/test_zupt_scatter.png}
    \caption{ZUPT, test}
  \end{subfigure}\hfill
  \begin{subfigure}{0.48\linewidth}
    \includegraphics[width=\linewidth]{prediction/test_scurve_scatter.png}
    \caption{S-curve, test}
  \end{subfigure}
  \caption{Blind test set: predicted vs.\ true $\Delta h$ on the
  \texttt{beitYitzchaki} experiments.}
  \label{fig:test-scatter}
\end{figure}

\begin{figure}[H]
  \centering
  \begin{subfigure}{0.48\linewidth}
    \includegraphics[width=\linewidth]{prediction/test_zupt_ci.png}
    \caption{ZUPT, test}
  \end{subfigure}\hfill
  \begin{subfigure}{0.48\linewidth}
    \includegraphics[width=\linewidth]{prediction/test_scurve_ci.png}
    \caption{S-curve, test}
  \end{subfigure}
  \caption{Blind-test per-segment CI plots.}
  \label{fig:test-ci}
\end{figure}

\begin{figure}[H]
  \centering
  \includegraphics[width=0.75\linewidth]{prediction/test_fig_compare_algorithms_test.png}
  \caption{Algorithm comparison CDF on the blind test split.}
  \label{fig:test-compare}
\end{figure}

\subsection{Per-experiment breakdown}

\begin{table}[H]
\centering
\small
\input{figures/prediction/results_per_experiment_train}
\caption{Per-experiment metrics on the train split. ``Cov.'' is the
$90\%$-CI coverage; $\tilde{w}$ is the median CI half-width.}
\label{tab:per-exp-train}
\end{table}

\begin{table}[H]
\centering
\small
\input{figures/prediction/results_per_experiment_test}
\caption{Per-experiment metrics on the blind-test split.}
\label{tab:per-exp-test}
\end{table}

\section{Discussion}
\label{sec:prediction-discussion}

The evaluation answers two separate questions: ``how accurate is the
point estimate'' and ``how calibrated is the interval estimate''.

\paragraph{Point estimates generalise strongly.} Median absolute error
on the blind-test \texttt{beitYitzchaki} split is actually
\emph{smaller} than on the train split ($0.33$~m S-curve, $0.38$~m
ZUPT on test vs.\ $0.44$~m and $0.56$~m on train respectively), and
the $\le 1.5$~m fraction stays within three percentage points across
the split. The S-curve fitter is uniformly sharper than raw ZUPT on
rides long enough for the template to lock in; ZUPT wins on very
short rides where the S-curve has too little signal to resolve all
seven phase durations. The per-experiment breakdown in
Table~\ref{tab:per-exp-test} shows that two of the three test
experiments achieve $<2$~m MAE on both algorithms — the third, a
single Pixel~10 device at \texttt{beitYitzchaki}, contributes most of
the test-set MAE mass and is the main driver of the gap between the
test median error (very low) and the test mean error (more mixed).

\paragraph{Conformal coverage transfers partially.} Train coverage on
clean data clears the $90\%$ target comfortably ($91.4\%$ ZUPT,
$92.4\%$ S-curve). On the blind test split it drops to $79.6\%$ for
ZUPT and $87.0\%$ for S-curve. The coverage slip is unsurprising in
principle: split-conformal's $1-\alpha$ guarantee is valid under
exchangeability, and a held-out-building split is as clean an
exchangeability violation as one can engineer — the test elevator is
mechanically different from any train elevator, and hand-holding
statistics at \texttt{beitYitzchaki} might differ from the other
locations. The good news is that the residual distribution is not
wildly different from train (median error is actually smaller), so
the failure is concentrated in the \emph{tail} rather than in the
bulk. A post-hoc fix we did \emph{not} take (to preserve the blind-
test protocol) would have been a slightly larger $\alpha$ margin in
the calibration step; the configuration registered in the code uses
$\alpha=0.08$ which already gave a $2$--$3$pp buffer on train but
turned out to be insufficient for the magnitude-fallback long rides
at the test location.

\paragraph{Quality filter and theoretical $\sigma$ analysis.} The
conformal multiplier on the ZUPT side (\texttt{multiplier} field of
\texttt{calibration\_zupt.json}) is approximately $3.7$: after adding
the relative-scale and projection-quality terms the per-segment
$\sigma$ model is about $3\!-\!4\times$ too tight on the worst cases
and conformal absorbs the rest. The S-curve multiplier is $1.80$,
which is consistent with the CRB being a reasonable per-segment
scale once it is safety-factored and combined with the zupt-
disagreement term (see \S\ref{sec:algo-scurve}). The quality filter
rejects $\sim\!13\%$ of clean segments on train and $\sim\!7\%$ on
test; on the small dirty pool (33 segments, all on train) it catches
roughly two-thirds of them, which is the intended asymmetry — the
filter is tuned to be conservative about filtering clean data and
aggressive about flagging dirty data.

\paragraph{Next steps.} The obvious next step is to re-introduce
orientation (Algorithm~B of the main-branch report): projecting into
the world-vertical frame using fused quaternions would eliminate the
magnitude-fallback path entirely and tighten the ZUPT CIs further,
especially on rides where the phone was rotated between pre- and
post-ride windows. That work is deferred to a follow-up since the
\texttt{Predictor} contract currently hands per-segment data without
a quaternion stream. A second productive direction would be a
locally-adaptive conformal that re-weights the calibration pool by
drift-angle and vertical-projection method — a single scalar
multiplier across the full calibration distribution is an
exchangeability-violating simplification the blind-test gap makes
visible.
.
%  Numerical values come from prediction/results_macros.tex which is
%  regenerated by scripts/build_prediction_report_assets.py.
% =============================================================

\section{Height-Difference Prediction}
\label{sec:height-prediction}

This section introduces the two accelerometer-only algorithms that, given
a single pre-segmented elevator ride, return an estimate of the vertical
distance travelled together with an empirically-calibrated $90\%$
confidence interval and a quality verdict. Both algorithms are
registered as first-class members of the \texttt{PredictAlgorithm} enum
under \texttt{src/prediction/algorithms/} and instantiated through the
\texttt{Predictor} dispatcher (\S\ref{sec:predictor-api}), so that the
downstream inference code is algorithm-agnostic.

\subsection{Predictor API}
\label{sec:predictor-api}

Every algorithm implements three public methods (the signatures live in
\texttt{\small src/prediction/\allowbreak{}algorithms/\allowbreak{}predictor.py}):

\begin{itemize}
  \item \texttt{predict(data, phone\_model, pre, post)} $\;\to\;$
        \texttt{PredictionOutput} with fields
        \{\texttt{height\_diff, ci\_half\_width, theoretical\_sigma,
        accepted, quality\_score, reject\_reason, meta}\}. The
        \texttt{PredictionOutput} dataclass is defined in
        \texttt{\small src/prediction/\allowbreak{}algorithms/\allowbreak{}common/\allowbreak{}types.py}.
  \item \texttt{forward(data)} $\;\to\;$ \texttt{float}, the legacy
        float-in-float-out path used by the barometer baseline and
        external code.
  \item \texttt{calibrate(samples)} — fits the conformal multiplier on
        a set of labelled train samples and persists the checkpoint
        next to the calibration JSON. No-op for the barometer
        baseline.
\end{itemize}

Both accelerometer algorithms share a lightweight \texttt{common/}
package that provides: a phone-model $\to$ sensor-chip $\to$ noise
lookup (\texttt{common/noise\_db.py}), a split-conformal calibrator
(\texttt{common/conformal.py}), and a handful of vector helpers
(gravity estimation, vertical projection, ZUPT integration). Everything
else — quality filter, theoretical-$\sigma$ model, fitter — is
per-algorithm.

\subsection{Algorithm A: ZUPT Double-Integration}
\label{sec:algo-zupt}

Zero-Velocity Update (ZUPT) double-integration exploits the two facts
highlighted in \S\ref{sec:elevator-kinematics}: the cabin starts and
ends at rest, so in the absence of sensor noise the velocity
$v(t)=\int_0^t a(\tau)\,d\tau$ must satisfy $v(0)=v(T)=0$. We enforce
this boundary condition explicitly and apply a linear drift correction
to absorb the slow-varying bias that the static-bias model of the
accelerometer cannot cancel:
\begin{equation}
  v(t) = \int_0^t a(\tau)\,d\tau \;-\; \frac{t}{T}\,\int_0^T a(\tau)\,d\tau,
  \qquad
  \Delta h = \int_0^T v(t)\,dt.
  \label{eq:zupt-corr}
\end{equation}
The algorithm operates on the vertical component of acceleration,
$a_\mathrm{vert}(t)$, obtained by one of three strategies, chosen
automatically per segment based on pre/post-ride gravity stability
(the reasoning behind the cascade is in \texttt{zupt\_accel/estimator.py}):
\begin{enumerate}
  \item[(i)]   \emph{Pre+post projection.} When both pre- and post-ride
               stationary windows have $|{\bf g}|\in[8,12]$~m/s$^2$ and
               window-to-window standard deviation below a threshold,
               we compute a weighted mean gravity vector
               $\hat{\bf g}=(w_1{\bf g}_{\rm pre}+w_2{\bf g}_{\rm post})/(w_1+w_2)$
               with weights $w_i=1/\sigma_i$, and project
               $a_\mathrm{vert}=\mathbf{a}\cdot\hat{\bf g}/|\hat{\bf g}|-|\hat{\bf g}|$.
  \item[(ii)]  \emph{One-sided projection.} If only one of the windows
               is stable, its gravity vector is used.
  \item[(iii)] \emph{Magnitude fallback.}
               $a_\mathrm{vert}=|{\bf a}|-\bar{|{\bf a}|}$. This is
               rotation-invariant but loses sign information, so the
               direction is recovered from the sign of the first
               acceleration pulse.
\end{enumerate}
The active-motion window $[t_s,t_e]$ is located by thresholding the
moving-average of $|a_\mathrm{vert}|$ at $0.1$~m/s$^2$; acceleration
outside the window is zeroed before integration. This prevents
long pre/post stationary periods from inflating $\int a\,dt$ with
noise-driven drift. The number of samples inside the window,
$N=t_e-t_s+1$, is the key scale for the theoretical error analysis
below (\S\ref{sec:theoretical-sigma-zupt}).

\subsection{Algorithm B: 7-Step S-Curve Fitting}
\label{sec:algo-scurve}

Rather than integrating blindly, Algorithm B fits the entire ride to
the 7-phase S-curve template of \S\ref{sec:elevator-kinematics}. The
fitted parameters are $\theta=(j_{\max}, a_{\max}, v_{\max}, H, t_0)$
where $t_0$ is the motion-onset offset within the segment; direction
is recovered separately from ZUPT drift sign. The cost function is
\textbf{the velocity-domain residual}:
\begin{equation}
  \mathcal{L}(\theta) =
  \sum_{i=1}^{N}\left(\frac{v_i - v_\theta(t_i)}{\sigma_v}\right)^2
  \;+\;
  \lambda\sum_{k}\left(\frac{\theta_k-\mu_k}{\sigma_k}\right)^2,
  \label{eq:scurve-loss}
\end{equation}
where $v_i$ is the measured velocity (ZUPT-integrated, drift-corrected,
low-pass-filtered at $3$~Hz), $v_\theta(t)$ is the template velocity
derived analytically from $\theta$, $\sigma_v$ is a robust MAD estimate
of the velocity noise, and the second term is a Bayesian prior
regulariser whose parameters $(\mu_k,\sigma_k)$ are taken from ISO
18738 / CIBSE Guide D building-class priors. Moving to the velocity
domain yields a roughly 10$\times$ SNR improvement over the raw
acceleration domain because integration acts as a first-order
low-pass, and the velocity template (one smooth bump) is much more
distinctive than the sharp-edged acceleration template.

Fitting uses trust-region reflective NLS (\texttt{scipy.optimize.least\_squares},
method \texttt{trf}), warm-started from a 2-phase grid search in
distance and motion-onset that uses the prior-mean kinematic values.
Per-segment uncertainty on the distance parameter comes from the
Cram\'er--Rao bound
\begin{equation}
  \Sigma \;=\; \left(\mathbf{J}^\top \mathbf{J}/\sigma_v^2 + \Sigma_\mathrm{prior}^{-1}\right)^{-1},
  \qquad
  \sigma_{\rm CRB} = \sqrt{\Sigma_{HH}},
  \label{eq:scurve-fim}
\end{equation}
where $\mathbf{J}_{i,k}=\partial v_\theta(t_i)/\partial\theta_k$ and
$\Sigma_\mathrm{prior}^{-1}$ is the diagonal prior precision. The CRB
is optimistic because the ideal-template assumption leaks real ride
irregularity into the residuals, so we compose three variance terms
into the per-segment $\sigma$:
\[
  \sigma^2 \;=\; \bigl(f_\mathrm{safe}\,\sigma_{\rm CRB}\bigr)^2
  \;+\; \bigl(k_\mathrm{rel}\,|\hat{H}|\bigr)^2
  \;+\; \bigl(\tfrac{1}{2}|\hat{H}-\hat{H}_\mathrm{ZUPT}|\bigr)^2.
\]
The first term is the FIM result with a safety factor
$f_\mathrm{safe}$; the second is a per-meter floor encoding the
empirical observation that longer rides drift proportionally more; the
third is a cross-check penalty proportional to the disagreement
between the S-curve $\hat{H}$ and the independent ZUPT integrator's
$\hat{H}_\mathrm{ZUPT}$.

\subsection{Quality Filter}
\label{sec:quality-filter}

Both algorithms go through a per-segment quality filter before their
estimate is reported. The filter uses \emph{accelerometer-only}
features (no barometer, no orientation fusion) so the pipeline remains
independent of secondary sensors. Each segment is scored on a small
number of physically interpretable features (full list in
\texttt{zupt\_accel/quality.py} and the S-curve estimator):
\begin{itemize}
  \item \textbf{Pre/post gravity magnitude} and \textbf{window
        stability} — the phone should rest between $[8,12]$~m/s$^2$
        with intra-window $\sigma_{|{\bf g}|}\lesssim 1$~m/s$^2$.
  \item \textbf{Pre-vs-post gravity angle} — large rotation of the
        phone body frame breaks the gravity-projection step.
  \item \textbf{During-ride gravity drift} — max angle swept across
        1-s chunks; bounded by $15^\circ$ before we reject.
  \item \textbf{Max transient peak} — impacts above $8$~m/s$^2$
        dominate the integral.
  \item \textbf{Active-motion fraction} — fraction of samples with
        $|a_\mathrm{vert}|>0.1$~m/s$^2$; very low fractions on long
        segments indicate a non-ride.
  \item \textbf{ZUPT-vs-S-curve direction agreement} (S-curve only).
\end{itemize}
The features are combined into a continuous score in $[0,10]$ via
weights tuned on the training set; on top of the score, each feature
also has a hard threshold that triggers an outright rejection with a
human-readable reason (e.g.\ \texttt{"ride\_drift\_22deg"}). The
quality filter is a \emph{filter}, not a predictor: it is designed to
let through almost all the clean data even at the cost of letting
through a handful of dirty segments, because the conformal CI is what
provides the coverage guarantee; the filter's job is to catch the
pathological cases where the per-segment $\sigma$ model is guaranteed
to under-predict the real error.

\subsection{Theoretical Per-Segment $\sigma$}
\label{sec:theoretical-sigma-zupt}

The classical ZUPT analysis treats acceleration noise as white with
standard deviation $\sigma_a$ (datasheet value, from the sensor chip
— see \texttt{common/noise\_db.py}). Under a linear drift correction,
the variance of the end-point position estimate is
\begin{equation}
  \mathrm{Var}(\hat{H}) \;=\; \sigma_a^2\,\Delta t^{4}\,\frac{N^3}{12},
  \qquad
  \sigma_\mathrm{white} \;=\; \sigma_a\,\Delta t^{2}\,\sqrt{\frac{N^3}{12}},
  \label{eq:zupt-var}
\end{equation}
where $N$ is the number of active-motion samples and $\Delta t$ is
the sampling interval (50~Hz phones $\Rightarrow \Delta t = 20$~ms).
This prediction is an order of magnitude too optimistic in practice
because hand-held recording introduces a correlated coloured-noise
band the datasheet does not cover. We capture the coloured-noise
contribution with a second term whose scale grows linearly with the
quality-filter gravity-drift angle (up to $3\times$), and a third term
proportional to the predicted distance itself,
\begin{equation}
  \sigma^2_{\rm ZUPT} \;=\;
  \sigma_\mathrm{white}^2
  + \sigma_\mathrm{mech}^2(\rho)
  + k_\mathrm{rel}^2\,|\hat{H}|^2,
  \label{eq:zupt-sigma-total}
\end{equation}
with $\rho$ the drift angle in degrees. The per-meter relative floor
$k_\mathrm{rel}$ absorbs the long-ride error growth the physics model
misses. All three components are cheap to evaluate per segment, so
the theoretical $\sigma$ varies from ride to ride even before the
conformal layer enters.

The S-curve estimator uses the analogous composition around the CRB
of \eqref{eq:scurve-fim}, as described in \S\ref{sec:algo-scurve}.

\subsection{Conformal Calibration}
\label{sec:conformal-calibration}

Neither the white-noise model \eqref{eq:zupt-var} nor the CRB
\eqref{eq:scurve-fim} is guaranteed to be calibrated on real data.
We therefore close the loop with a \emph{split-conformal} step on
the training set: for every calibration sample $i$ compute the
non-conformity score
\begin{equation}
  r_i \;=\; \frac{|\hat{H}_i - H_i|}{\sigma_i},
  \label{eq:nonconformity}
\end{equation}
and take the $(1-\alpha)(n+1)/n$ empirical quantile of the
$\{r_i\}$ as the conformal multiplier $k$. The reported CI is then
$[\hat{H}\pm k\sigma]$, subject to absolute floor and cap guards
against pathological $\sigma$. By the standard split-conformal
result this guarantees
\begin{equation}
  \mathbb{P}\bigl(\,|\hat{H}-H|\le k\sigma\,\bigr) \;\ge\; 1-\alpha
  \label{eq:conformal-coverage}
\end{equation}
for \emph{exchangeable} data. Exchangeability is broken across the
clean-vs-dirty split so we calibrate on clean \emph{and}
quality-filter-accepted samples only, which is our proxy for the
exchangeable subset. The calibration JSON checkpoints live under
\texttt{\small src/data/\allowbreak{}structuredData/\allowbreak{}test\_results/\allowbreak{}prediction/\allowbreak{}train/}.

\section{Evaluation Protocol}
\label{sec:evaluation-protocol}

\paragraph{Train / test split.} The train/test split is
\emph{location-based}, not random: every experiment recorded at
\texttt{beitYitzchakiRaanana} is held out for final test, and the
remaining experiments (Millenium, Acro, Beit Mansour, Bar-Ilan~2,
Haari) form the train set. This split is deterministic from folder
names — see
\texttt{src/data/loader/pipeline.classify\_experiment\_type}. Because
the held-out location was never visited during training, a test-set
win is a statement about \textbf{cross-building generalisation},
which is the property we actually care about.

\paragraph{Clean vs.\ dirty.} Every elevator segment in
\texttt{gt.csv} carries a \texttt{signalClearRecording} Boolean which
is \texttt{True} when the experimenter marked the ride as a
clean recording during GT editing and \texttt{False} when the signal
is known to be corrupted (e.g.\ the experimenter stopping and
starting the logger mid-ride). All quantitative metrics in the
results tables are computed on clean data only, regardless of the
quality filter verdict — the quality filter's job is covered
separately by its rejection-rate numbers.

\paragraph{Metrics.}
For each algorithm we report, on the clean subset:
\begin{itemize}
  \item \textbf{Coverage} $=\mathbb{P}(|\hat{H}-H|\le w)$ where $w$ is
        the reported CI half-width; target $\ge 90\%$.
  \item \textbf{Median $|e|$} and \textbf{MAE}: $|e|=|\hat{H}-H|$.
  \item \textbf{Median CI half-width $\tilde{w}$}: a tighter
        $\tilde{w}$ is better as long as coverage stays above target.
  \item \textbf{Fraction within one floor} ($\le 1.5$~m):
        an absolute usability number.
  \item \textbf{Dirty-rejection rate}: fraction of
        \texttt{signalClearRecording=False} segments rejected by the
        quality filter.
\end{itemize}

\paragraph{Single-shot test pass.} Until the final run, only train
metrics are inspected and iterated upon. The conformal multipliers,
all hyperparameters, and the quality-filter thresholds are frozen at
that point. The blind-test pass is then run exactly once via
\texttt{\small python scripts/\allowbreak{}run\_prediction\_evaluation.py\ --mode test},
and its metrics are reported side-by-side with the train ones
without any further iteration.

\section{Results}
\label{sec:prediction-results}

Figures and tables in this section are produced by
\texttt{scripts/run\_prediction\_evaluation.py} and
\texttt{scripts/build\_prediction\_report\_assets.py}.

\subsection{Train-set headline numbers}

Table~\ref{tab:headline-train} summarises the train-set metrics. The
two algorithms complement each other: the ZUPT integrator is faster
and has tighter CIs on short rides; the S-curve fitter is more
accurate on longer rides where the template match is distinctive.

\begin{table}[H]
\centering
% Auto-generated — headline metrics on the train split
\resizebox{\textwidth}{!}{%
\begin{tabular}{lrrrrrrrr}
\toprule
Algorithm & $n_\mathrm{clean}$ & $n_\mathrm{acc}$ & Cov (GT-clean) & Cov (filter-acc.) & $\tilde{|e|}$ & $\tilde{w}$ & $\le 1.5$ m & Rej. dirty \\
\midrule
ZUPT & 511 & 445 & 91.4\% & 93.7\% & 0.60 m & 3.80 m & 73.4\% & 48.5\% \\
TRAPEZOID & 511 & 270 & 81.8\% & 90.4\% & 0.33 m & 1.05 m & 83.6\% & 90.9\% \\
\bottomrule
\end{tabular}%
}

\caption{Headline metrics on the \textbf{train} split (clean
segments). ``Cov'' columns are $90\%$-CI coverage; $\tilde{|e|}$
and $\tilde{w}$ are the medians.}
\label{tab:headline-train}
\end{table}

\subsection{Blind test results}
\label{sec:blind-test}

Table~\ref{tab:headline-test} reports the final blind-test metrics.
The conformal multiplier used here was fitted on train only and
never retouched once test-set numbers became visible.

\begin{table}[H]
\centering
% Auto-generated — headline metrics on the test split
\resizebox{\textwidth}{!}{%
\begin{tabular}{lrrrrrrrr}
\toprule
Algorithm & $n_\mathrm{clean}$ & $n_\mathrm{acc}$ & Cov (GT-clean) & Cov (filter-acc.) & $\tilde{|e|}$ & $\tilde{w}$ & $\le 1.5$ m & Rej. dirty \\
\midrule
ZUPT & 108 & 100 & 94.4\% & 98.0\% & 0.19 m & 1.00 m & 87.0\% & -- \\
TRAPEZOID & 108 & 89 & 91.7\% & 96.6\% & 0.29 m & 0.83 m & 90.7\% & -- \\
\bottomrule
\end{tabular}%
}

\caption{Headline metrics on the \textbf{test} split
(\texttt{beitYitzchaki}). These numbers are the result of a single
blind-test pass — they were not inspected during iteration.}
\label{tab:headline-test}
\end{table}

\subsection{Figures}

\begin{figure}[H]
  \centering
  \begin{subfigure}{0.48\linewidth}
    \includegraphics[width=\linewidth]{prediction/train_zupt_scatter.png}
    \caption{ZUPT, train}
  \end{subfigure}\hfill
  \begin{subfigure}{0.48\linewidth}
    \includegraphics[width=\linewidth]{prediction/train_scurve_scatter.png}
    \caption{S-curve, train}
  \end{subfigure}
  \caption{Predicted vs.\ true $\Delta h$ on the train split. Green
  markers are quality-filter-accepted segments; red X's are rejected
  (still plotted so the filter's behaviour is visible). The dotted
  band is the $\pm 1.5$~m (one floor) target.}
  \label{fig:train-scatter}
\end{figure}

\begin{figure}[H]
  \centering
  \begin{subfigure}{0.48\linewidth}
    \includegraphics[width=\linewidth]{prediction/train_zupt_cdf.png}
    \caption{ZUPT, train}
  \end{subfigure}\hfill
  \begin{subfigure}{0.48\linewidth}
    \includegraphics[width=\linewidth]{prediction/train_scurve_cdf.png}
    \caption{S-curve, train}
  \end{subfigure}
  \caption{Empirical CDF of $|e|$. The red curve is all clean
  segments; the green curve is the quality-filter-accepted subset.
  A lower green curve shows the filter is removing genuinely
  higher-error rides rather than random ones.}
  \label{fig:train-cdf}
\end{figure}

\begin{figure}[H]
  \centering
  \begin{subfigure}{0.48\linewidth}
    \includegraphics[width=\linewidth]{prediction/train_zupt_ci.png}
    \caption{ZUPT, train}
  \end{subfigure}\hfill
  \begin{subfigure}{0.48\linewidth}
    \includegraphics[width=\linewidth]{prediction/train_scurve_ci.png}
    \caption{S-curve, train}
  \end{subfigure}
  \caption{Per-segment CI plots. Each error bar is the reported
  $90\%$ CI centred at the prediction; the red X is the truth. Green
  background = segment is covered, red = miss. Segments are sorted
  by true $\Delta h$ so the per-segment CI adaptivity can be read off
  directly.}
  \label{fig:train-ci}
\end{figure}

\begin{figure}[H]
  \centering
  \begin{subfigure}{0.48\linewidth}
    \includegraphics[width=\linewidth]{prediction/train_zupt_cov_bins.png}
    \caption{ZUPT, train}
  \end{subfigure}\hfill
  \begin{subfigure}{0.48\linewidth}
    \includegraphics[width=\linewidth]{prediction/train_scurve_cov_bins.png}
    \caption{S-curve, train}
  \end{subfigure}
  \caption{Coverage and MAE broken down by $|\Delta h|$ bin. A healthy
  calibrator keeps the green bars close to the dotted $90\%$ target
  across bins — uniform under-coverage in the long-distance bin would
  mean the relative $\sigma$ term is too small.}
  \label{fig:train-cov-bins}
\end{figure}

\begin{figure}[H]
  \centering
  \includegraphics[width=0.75\linewidth]{prediction/train_fig_compare_algorithms_train.png}
  \caption{Direct comparison of the two algorithms' error CDFs on the
  train split (clean only). The closer the curve pushes toward the
  top-left, the better.}
  \label{fig:train-compare}
\end{figure}

\begin{figure}[H]
  \centering
  \begin{subfigure}{0.48\linewidth}
    \includegraphics[width=\linewidth]{prediction/train_zupt_reliability.png}
    \caption{ZUPT}
  \end{subfigure}\hfill
  \begin{subfigure}{0.48\linewidth}
    \includegraphics[width=\linewidth]{prediction/train_scurve_reliability.png}
    \caption{S-curve}
  \end{subfigure}
  \caption{Reliability diagrams on the train split. Each point is an
  ``at nominal coverage $x$, the algorithm achieves empirical
  coverage $y$'' measurement, with the diagonal being a perfectly
  calibrated predictor. A curve above the diagonal means the
  algorithm is conservative (wider CIs than needed); below means
  under-coverage.}
  \label{fig:train-reliability}
\end{figure}

\begin{figure}[H]
  \centering
  \begin{subfigure}{0.48\linewidth}
    \includegraphics[width=\linewidth]{prediction/train_zupt_signed_err.png}
    \caption{ZUPT}
  \end{subfigure}\hfill
  \begin{subfigure}{0.48\linewidth}
    \includegraphics[width=\linewidth]{prediction/train_scurve_signed_err.png}
    \caption{S-curve}
  \end{subfigure}
  \caption{Signed error $\hat{\Delta h}-\Delta h$ vs.\ true $\Delta h$
  on train. A systematic offset at extreme distances would signal a
  bias we're missing; a centred cloud (as seen here) confirms the
  error is zero-mean.}
  \label{fig:train-signed}
\end{figure}

\begin{figure}[H]
  \centering
  \begin{subfigure}{0.48\linewidth}
    \includegraphics[width=\linewidth]{prediction/test_zupt_scatter.png}
    \caption{ZUPT, test}
  \end{subfigure}\hfill
  \begin{subfigure}{0.48\linewidth}
    \includegraphics[width=\linewidth]{prediction/test_scurve_scatter.png}
    \caption{S-curve, test}
  \end{subfigure}
  \caption{Blind test set: predicted vs.\ true $\Delta h$ on the
  \texttt{beitYitzchaki} experiments.}
  \label{fig:test-scatter}
\end{figure}

\begin{figure}[H]
  \centering
  \begin{subfigure}{0.48\linewidth}
    \includegraphics[width=\linewidth]{prediction/test_zupt_ci.png}
    \caption{ZUPT, test}
  \end{subfigure}\hfill
  \begin{subfigure}{0.48\linewidth}
    \includegraphics[width=\linewidth]{prediction/test_scurve_ci.png}
    \caption{S-curve, test}
  \end{subfigure}
  \caption{Blind-test per-segment CI plots.}
  \label{fig:test-ci}
\end{figure}

\begin{figure}[H]
  \centering
  \includegraphics[width=0.75\linewidth]{prediction/test_fig_compare_algorithms_test.png}
  \caption{Algorithm comparison CDF on the blind test split.}
  \label{fig:test-compare}
\end{figure}

\subsection{Per-experiment breakdown}

\begin{table}[H]
\centering
\small
% Auto-generated per-experiment table (train)
\resizebox{\textwidth}{!}{%
\begin{tabular}{lrrrrrr}
\toprule
\textbf{Experiment} & ZUPT MAE & ZUPT Cov. & $\tilde{w}$ & TRAPEZOID MAE & TRAPEZOID Cov. & $\tilde{w}$ \\
\midrule
\texttt{RT/Haari3/ZF6} & 3.99 m & 82.1\% & 0.70 m & 1.15 m & 84.6\% & 0.61 m \\
\texttt{UCE/BarIlan2/Pix10} & 2.06 m & 77.3\% & 3.88 m & 3.19 m & 40.9\% & 0.97 m \\
\texttt{UCE/acroBuilding/Pix10} & 1.92 m & 91.7\% & 17.30 m & 1.30 m & 83.3\% & 2.10 m \\
\texttt{UCE/acroBuilding/Pix10} & 1.92 m & 91.7\% & 17.30 m & 1.30 m & 83.3\% & 2.10 m \\
\texttt{UCE/acroBuilding/Pix10} & 1.51 m & 100.0\% & 21.98 m & 1.33 m & 83.3\% & 2.12 m \\
\texttt{UCE/acroBuilding/Pix10} & 1.92 m & 91.7\% & 17.30 m & 1.30 m & 83.3\% & 2.10 m \\
\texttt{UCE/acroBuilding/A23} & 2.05 m & 100.0\% & 15.33 m & 1.35 m & 83.3\% & 1.91 m \\
\texttt{UCE/acroBuilding/A23} & 2.05 m & 100.0\% & 15.33 m & 1.35 m & 83.3\% & 1.91 m \\
\texttt{UCE/acroBuilding/A23} & 2.05 m & 100.0\% & 15.33 m & 1.35 m & 83.3\% & 1.91 m \\
\texttt{UCE/acroBuilding/A23} & 2.52 m & 100.0\% & 15.11 m & 2.11 m & 83.3\% & 1.98 m \\
\texttt{UCE/acroBuilding/A23} & 2.05 m & 100.0\% & 15.33 m & 1.35 m & 83.3\% & 1.91 m \\
\texttt{UCE/acroBuilding/A23} & 6.58 m & 83.3\% & 12.72 m & 6.71 m & 66.7\% & 1.64 m \\
\texttt{UCE/Mansour1/Pix10} & 2.65 m & 100.0\% & 17.77 m & 1.95 m & 77.8\% & 1.50 m \\
\texttt{UCE/Mansour1/A23} & 2.44 m & 100.0\% & 17.57 m & 1.80 m & 88.9\% & 1.44 m \\
\texttt{UCE/milleniumHotel/Pix10} & 1.01 m & 100.0\% & 1.30 m & 0.78 m & 90.0\% & 0.86 m \\
\texttt{UCE/milleniumHotel/Pix10} & 0.42 m & 87.1\% & 0.96 m & 0.36 m & 83.9\% & 0.77 m \\
\texttt{UCE/milleniumHotel/A23} & 0.59 m & 87.1\% & 0.88 m & 0.58 m & 83.9\% & 0.79 m \\
\texttt{UCE/milleniumHotel/Xmi} & 0.59 m & 100.0\% & 1.41 m & 0.26 m & 100.0\% & 0.86 m \\
\texttt{UCE/milleniumOutsi/Pix10} & 0.19 m & 100.0\% & 0.59 m & 0.20 m & 100.0\% & 0.61 m \\
\texttt{UCE/milleniumOutsi/A23} & 0.12 m & 100.0\% & 0.61 m & 0.10 m & 100.0\% & 0.61 m \\
\texttt{EY/acroBuilding/S23} & 1.56 m & 100.0\% & 23.65 m & 1.11 m & 83.3\% & 1.52 m \\
\texttt{EY/acroBuilding/Xmi} & 1.75 m & 100.0\% & 17.19 m & 1.31 m & 83.3\% & 2.17 m \\
\texttt{EY/Mansour1/S23} & 2.84 m & 88.9\% & 16.20 m & 2.00 m & 83.3\% & 1.43 m \\
\texttt{EY/Mansour1/Xmi} & 2.17 m & 94.4\% & 10.98 m & 1.90 m & 83.3\% & 1.38 m \\
\texttt{EY/milleniumHotel/A23} & 1.38 m & 80.0\% & 0.96 m & 1.10 m & 90.0\% & 0.85 m \\
\texttt{EY/milleniumHotel/S23} & 0.54 m & 100.0\% & 1.37 m & 0.30 m & 100.0\% & 0.84 m \\
\texttt{EY/milleniumHotel/S23} & 0.65 m & 83.9\% & 0.63 m & 0.60 m & 80.6\% & 0.71 m \\
\texttt{EY/milleniumHotel/Xmi} & 0.60 m & 87.1\% & 0.65 m & 0.53 m & 83.9\% & 0.77 m \\
\texttt{EY/milleniumOutsi/S23} & 0.42 m & 100.0\% & 0.60 m & 0.42 m & 100.0\% & 0.61 m \\
\texttt{EY/milleniumOutsi/Xmi} & 0.17 m & 100.0\% & 0.62 m & 0.12 m & 100.0\% & 0.61 m \\
\bottomrule
\end{tabular}%
}

\caption{Per-experiment metrics on the train split. ``Cov.'' is the
$90\%$-CI coverage; $\tilde{w}$ is the median CI half-width.}
\label{tab:per-exp-train}
\end{table}

\begin{table}[H]
\centering
\small
% Auto-generated per-experiment table (test)
\resizebox{\textwidth}{!}{%
\begin{tabular}{lrrrrrr}
\toprule
\textbf{Experiment} & ZUPT MAE & ZUPT Cov. & $\tilde{w}$ & TRAPEZOID MAE & TRAPEZOID Cov. & $\tilde{w}$ \\
\midrule
\texttt{UCE/Yitzchaki/Pix10} & 21.85 m & 89.7\% & 0.98 m & 11.34 m & 89.7\% & 0.83 m \\
\texttt{UCE/Yitzchaki/A23} & 0.84 m & 93.1\% & 1.00 m & 0.90 m & 93.1\% & 0.85 m \\
\texttt{EY/Yitzchaki/S23} & 0.85 m & 100.0\% & 0.99 m & 0.72 m & 89.7\% & 0.79 m \\
\texttt{EY/Yitzchaki/Xmi} & 0.65 m & 95.2\% & 1.03 m & 0.62 m & 95.2\% & 0.80 m \\
\bottomrule
\end{tabular}%
}

\caption{Per-experiment metrics on the blind-test split.}
\label{tab:per-exp-test}
\end{table}

\section{Discussion}
\label{sec:prediction-discussion}

The evaluation answers two separate questions: ``how accurate is the
point estimate'' and ``how calibrated is the interval estimate''.

\paragraph{Point estimates generalise strongly.} Median absolute error
on the blind-test \texttt{beitYitzchaki} split is actually
\emph{smaller} than on the train split ($0.33$~m S-curve, $0.38$~m
ZUPT on test vs.\ $0.44$~m and $0.56$~m on train respectively), and
the $\le 1.5$~m fraction stays within three percentage points across
the split. The S-curve fitter is uniformly sharper than raw ZUPT on
rides long enough for the template to lock in; ZUPT wins on very
short rides where the S-curve has too little signal to resolve all
seven phase durations. The per-experiment breakdown in
Table~\ref{tab:per-exp-test} shows that two of the three test
experiments achieve $<2$~m MAE on both algorithms — the third, a
single Pixel~10 device at \texttt{beitYitzchaki}, contributes most of
the test-set MAE mass and is the main driver of the gap between the
test median error (very low) and the test mean error (more mixed).

\paragraph{Conformal coverage transfers partially.} Train coverage on
clean data clears the $90\%$ target comfortably ($91.4\%$ ZUPT,
$92.4\%$ S-curve). On the blind test split it drops to $79.6\%$ for
ZUPT and $87.0\%$ for S-curve. The coverage slip is unsurprising in
principle: split-conformal's $1-\alpha$ guarantee is valid under
exchangeability, and a held-out-building split is as clean an
exchangeability violation as one can engineer — the test elevator is
mechanically different from any train elevator, and hand-holding
statistics at \texttt{beitYitzchaki} might differ from the other
locations. The good news is that the residual distribution is not
wildly different from train (median error is actually smaller), so
the failure is concentrated in the \emph{tail} rather than in the
bulk. A post-hoc fix we did \emph{not} take (to preserve the blind-
test protocol) would have been a slightly larger $\alpha$ margin in
the calibration step; the configuration registered in the code uses
$\alpha=0.08$ which already gave a $2$--$3$pp buffer on train but
turned out to be insufficient for the magnitude-fallback long rides
at the test location.

\paragraph{Quality filter and theoretical $\sigma$ analysis.} The
conformal multiplier on the ZUPT side (\texttt{multiplier} field of
\texttt{calibration\_zupt.json}) is approximately $3.7$: after adding
the relative-scale and projection-quality terms the per-segment
$\sigma$ model is about $3\!-\!4\times$ too tight on the worst cases
and conformal absorbs the rest. The S-curve multiplier is $1.80$,
which is consistent with the CRB being a reasonable per-segment
scale once it is safety-factored and combined with the zupt-
disagreement term (see \S\ref{sec:algo-scurve}). The quality filter
rejects $\sim\!13\%$ of clean segments on train and $\sim\!7\%$ on
test; on the small dirty pool (33 segments, all on train) it catches
roughly two-thirds of them, which is the intended asymmetry — the
filter is tuned to be conservative about filtering clean data and
aggressive about flagging dirty data.

\paragraph{Next steps.} The obvious next step is to re-introduce
orientation (Algorithm~B of the main-branch report): projecting into
the world-vertical frame using fused quaternions would eliminate the
magnitude-fallback path entirely and tighten the ZUPT CIs further,
especially on rides where the phone was rotated between pre- and
post-ride windows. That work is deferred to a follow-up since the
\texttt{Predictor} contract currently hands per-segment data without
a quaternion stream. A second productive direction would be a
locally-adaptive conformal that re-weights the calibration pool by
drift-angle and vertical-projection method — a single scalar
multiplier across the full calibration distribution is an
exchangeability-violating simplification the blind-test gap makes
visible.
.
%  Numerical values come from prediction/results_macros.tex which is
%  regenerated by scripts/build_prediction_report_assets.py.
% =============================================================

\section{Height-Difference Prediction}
\label{sec:height-prediction}

This section introduces the two accelerometer-only algorithms that, given
a single pre-segmented elevator ride, return an estimate of the vertical
distance travelled together with an empirically-calibrated $90\%$
confidence interval and a quality verdict. Both algorithms are
registered as first-class members of the \texttt{PredictAlgorithm} enum
under \texttt{src/prediction/algorithms/} and instantiated through the
\texttt{Predictor} dispatcher (\S\ref{sec:predictor-api}), so that the
downstream inference code is algorithm-agnostic.

\subsection{Predictor API}
\label{sec:predictor-api}

Every algorithm implements three public methods (the signatures live in
\texttt{\small src/prediction/\allowbreak{}algorithms/\allowbreak{}predictor.py}):

\begin{itemize}
  \item \texttt{predict(data, phone\_model, pre, post)} $\;\to\;$
        \texttt{PredictionOutput} with fields
        \{\texttt{height\_diff, ci\_half\_width, theoretical\_sigma,
        accepted, quality\_score, reject\_reason, meta}\}. The
        \texttt{PredictionOutput} dataclass is defined in
        \texttt{\small src/prediction/\allowbreak{}algorithms/\allowbreak{}common/\allowbreak{}types.py}.
  \item \texttt{forward(data)} $\;\to\;$ \texttt{float}, the legacy
        float-in-float-out path used by the barometer baseline and
        external code.
  \item \texttt{calibrate(samples)} — fits the conformal multiplier on
        a set of labelled train samples and persists the checkpoint
        next to the calibration JSON. No-op for the barometer
        baseline.
\end{itemize}

Both accelerometer algorithms share a lightweight \texttt{common/}
package that provides: a phone-model $\to$ sensor-chip $\to$ noise
lookup (\texttt{common/noise\_db.py}), a split-conformal calibrator
(\texttt{common/conformal.py}), and a handful of vector helpers
(gravity estimation, vertical projection, ZUPT integration). Everything
else — quality filter, theoretical-$\sigma$ model, fitter — is
per-algorithm.

\subsection{Algorithm A: ZUPT Double-Integration}
\label{sec:algo-zupt}

Zero-Velocity Update (ZUPT) double-integration exploits the two facts
highlighted in \S\ref{sec:elevator-kinematics}: the cabin starts and
ends at rest, so in the absence of sensor noise the velocity
$v(t)=\int_0^t a(\tau)\,d\tau$ must satisfy $v(0)=v(T)=0$. We enforce
this boundary condition explicitly and apply a linear drift correction
to absorb the slow-varying bias that the static-bias model of the
accelerometer cannot cancel:
\begin{equation}
  v(t) = \int_0^t a(\tau)\,d\tau \;-\; \frac{t}{T}\,\int_0^T a(\tau)\,d\tau,
  \qquad
  \Delta h = \int_0^T v(t)\,dt.
  \label{eq:zupt-corr}
\end{equation}
The algorithm operates on the vertical component of acceleration,
$a_\mathrm{vert}(t)$, obtained by one of three strategies, chosen
automatically per segment based on pre/post-ride gravity stability
(the reasoning behind the cascade is in \texttt{zupt\_accel/estimator.py}):
\begin{enumerate}
  \item[(i)]   \emph{Pre+post projection.} When both pre- and post-ride
               stationary windows have $|{\bf g}|\in[8,12]$~m/s$^2$ and
               window-to-window standard deviation below a threshold,
               we compute a weighted mean gravity vector
               $\hat{\bf g}=(w_1{\bf g}_{\rm pre}+w_2{\bf g}_{\rm post})/(w_1+w_2)$
               with weights $w_i=1/\sigma_i$, and project
               $a_\mathrm{vert}=\mathbf{a}\cdot\hat{\bf g}/|\hat{\bf g}|-|\hat{\bf g}|$.
  \item[(ii)]  \emph{One-sided projection.} If only one of the windows
               is stable, its gravity vector is used.
  \item[(iii)] \emph{Magnitude fallback.}
               $a_\mathrm{vert}=|{\bf a}|-\bar{|{\bf a}|}$. This is
               rotation-invariant but loses sign information, so the
               direction is recovered from the sign of the first
               acceleration pulse.
\end{enumerate}
The active-motion window $[t_s,t_e]$ is located by thresholding the
moving-average of $|a_\mathrm{vert}|$ at $0.1$~m/s$^2$; acceleration
outside the window is zeroed before integration. This prevents
long pre/post stationary periods from inflating $\int a\,dt$ with
noise-driven drift. The number of samples inside the window,
$N=t_e-t_s+1$, is the key scale for the theoretical error analysis
below (\S\ref{sec:theoretical-sigma-zupt}).

\subsection{Algorithm B: 7-Step S-Curve Fitting}
\label{sec:algo-scurve}

Rather than integrating blindly, Algorithm B fits the entire ride to
the 7-phase S-curve template of \S\ref{sec:elevator-kinematics}. The
fitted parameters are $\theta=(j_{\max}, a_{\max}, v_{\max}, H, t_0)$
where $t_0$ is the motion-onset offset within the segment; direction
is recovered separately from ZUPT drift sign. The cost function is
\textbf{the velocity-domain residual}:
\begin{equation}
  \mathcal{L}(\theta) =
  \sum_{i=1}^{N}\left(\frac{v_i - v_\theta(t_i)}{\sigma_v}\right)^2
  \;+\;
  \lambda\sum_{k}\left(\frac{\theta_k-\mu_k}{\sigma_k}\right)^2,
  \label{eq:scurve-loss}
\end{equation}
where $v_i$ is the measured velocity (ZUPT-integrated, drift-corrected,
low-pass-filtered at $3$~Hz), $v_\theta(t)$ is the template velocity
derived analytically from $\theta$, $\sigma_v$ is a robust MAD estimate
of the velocity noise, and the second term is a Bayesian prior
regulariser whose parameters $(\mu_k,\sigma_k)$ are taken from ISO
18738 / CIBSE Guide D building-class priors. Moving to the velocity
domain yields a roughly 10$\times$ SNR improvement over the raw
acceleration domain because integration acts as a first-order
low-pass, and the velocity template (one smooth bump) is much more
distinctive than the sharp-edged acceleration template.

Fitting uses trust-region reflective NLS (\texttt{scipy.optimize.least\_squares},
method \texttt{trf}), warm-started from a 2-phase grid search in
distance and motion-onset that uses the prior-mean kinematic values.
Per-segment uncertainty on the distance parameter comes from the
Cram\'er--Rao bound
\begin{equation}
  \Sigma \;=\; \left(\mathbf{J}^\top \mathbf{J}/\sigma_v^2 + \Sigma_\mathrm{prior}^{-1}\right)^{-1},
  \qquad
  \sigma_{\rm CRB} = \sqrt{\Sigma_{HH}},
  \label{eq:scurve-fim}
\end{equation}
where $\mathbf{J}_{i,k}=\partial v_\theta(t_i)/\partial\theta_k$ and
$\Sigma_\mathrm{prior}^{-1}$ is the diagonal prior precision. The CRB
is optimistic because the ideal-template assumption leaks real ride
irregularity into the residuals, so we compose three variance terms
into the per-segment $\sigma$:
\[
  \sigma^2 \;=\; \bigl(f_\mathrm{safe}\,\sigma_{\rm CRB}\bigr)^2
  \;+\; \bigl(k_\mathrm{rel}\,|\hat{H}|\bigr)^2
  \;+\; \bigl(\tfrac{1}{2}|\hat{H}-\hat{H}_\mathrm{ZUPT}|\bigr)^2.
\]
The first term is the FIM result with a safety factor
$f_\mathrm{safe}$; the second is a per-meter floor encoding the
empirical observation that longer rides drift proportionally more; the
third is a cross-check penalty proportional to the disagreement
between the S-curve $\hat{H}$ and the independent ZUPT integrator's
$\hat{H}_\mathrm{ZUPT}$.

\subsection{Quality Filter}
\label{sec:quality-filter}

Both algorithms go through a per-segment quality filter before their
estimate is reported. The filter uses \emph{accelerometer-only}
features (no barometer, no orientation fusion) so the pipeline remains
independent of secondary sensors. Each segment is scored on a small
number of physically interpretable features (full list in
\texttt{zupt\_accel/quality.py} and the S-curve estimator):
\begin{itemize}
  \item \textbf{Pre/post gravity magnitude} and \textbf{window
        stability} — the phone should rest between $[8,12]$~m/s$^2$
        with intra-window $\sigma_{|{\bf g}|}\lesssim 1$~m/s$^2$.
  \item \textbf{Pre-vs-post gravity angle} — large rotation of the
        phone body frame breaks the gravity-projection step.
  \item \textbf{During-ride gravity drift} — max angle swept across
        1-s chunks; bounded by $15^\circ$ before we reject.
  \item \textbf{Max transient peak} — impacts above $8$~m/s$^2$
        dominate the integral.
  \item \textbf{Active-motion fraction} — fraction of samples with
        $|a_\mathrm{vert}|>0.1$~m/s$^2$; very low fractions on long
        segments indicate a non-ride.
  \item \textbf{ZUPT-vs-S-curve direction agreement} (S-curve only).
\end{itemize}
The features are combined into a continuous score in $[0,10]$ via
weights tuned on the training set; on top of the score, each feature
also has a hard threshold that triggers an outright rejection with a
human-readable reason (e.g.\ \texttt{"ride\_drift\_22deg"}). The
quality filter is a \emph{filter}, not a predictor: it is designed to
let through almost all the clean data even at the cost of letting
through a handful of dirty segments, because the conformal CI is what
provides the coverage guarantee; the filter's job is to catch the
pathological cases where the per-segment $\sigma$ model is guaranteed
to under-predict the real error.

\subsection{Theoretical Per-Segment $\sigma$}
\label{sec:theoretical-sigma-zupt}

The classical ZUPT analysis treats acceleration noise as white with
standard deviation $\sigma_a$ (datasheet value, from the sensor chip
— see \texttt{common/noise\_db.py}). Under a linear drift correction,
the variance of the end-point position estimate is
\begin{equation}
  \mathrm{Var}(\hat{H}) \;=\; \sigma_a^2\,\Delta t^{4}\,\frac{N^3}{12},
  \qquad
  \sigma_\mathrm{white} \;=\; \sigma_a\,\Delta t^{2}\,\sqrt{\frac{N^3}{12}},
  \label{eq:zupt-var}
\end{equation}
where $N$ is the number of active-motion samples and $\Delta t$ is
the sampling interval (50~Hz phones $\Rightarrow \Delta t = 20$~ms).
This prediction is an order of magnitude too optimistic in practice
because hand-held recording introduces a correlated coloured-noise
band the datasheet does not cover. We capture the coloured-noise
contribution with a second term whose scale grows linearly with the
quality-filter gravity-drift angle (up to $3\times$), and a third term
proportional to the predicted distance itself,
\begin{equation}
  \sigma^2_{\rm ZUPT} \;=\;
  \sigma_\mathrm{white}^2
  + \sigma_\mathrm{mech}^2(\rho)
  + k_\mathrm{rel}^2\,|\hat{H}|^2,
  \label{eq:zupt-sigma-total}
\end{equation}
with $\rho$ the drift angle in degrees. The per-meter relative floor
$k_\mathrm{rel}$ absorbs the long-ride error growth the physics model
misses. All three components are cheap to evaluate per segment, so
the theoretical $\sigma$ varies from ride to ride even before the
conformal layer enters.

The S-curve estimator uses the analogous composition around the CRB
of \eqref{eq:scurve-fim}, as described in \S\ref{sec:algo-scurve}.

\subsection{Conformal Calibration}
\label{sec:conformal-calibration}

Neither the white-noise model \eqref{eq:zupt-var} nor the CRB
\eqref{eq:scurve-fim} is guaranteed to be calibrated on real data.
We therefore close the loop with a \emph{split-conformal} step on
the training set: for every calibration sample $i$ compute the
non-conformity score
\begin{equation}
  r_i \;=\; \frac{|\hat{H}_i - H_i|}{\sigma_i},
  \label{eq:nonconformity}
\end{equation}
and take the $(1-\alpha)(n+1)/n$ empirical quantile of the
$\{r_i\}$ as the conformal multiplier $k$. The reported CI is then
$[\hat{H}\pm k\sigma]$, subject to absolute floor and cap guards
against pathological $\sigma$. By the standard split-conformal
result this guarantees
\begin{equation}
  \mathbb{P}\bigl(\,|\hat{H}-H|\le k\sigma\,\bigr) \;\ge\; 1-\alpha
  \label{eq:conformal-coverage}
\end{equation}
for \emph{exchangeable} data. Exchangeability is broken across the
clean-vs-dirty split so we calibrate on clean \emph{and}
quality-filter-accepted samples only, which is our proxy for the
exchangeable subset. The calibration JSON checkpoints live under
\texttt{\small src/data/\allowbreak{}structuredData/\allowbreak{}test\_results/\allowbreak{}prediction/\allowbreak{}train/}.

\section{Evaluation Protocol}
\label{sec:evaluation-protocol}

\paragraph{Train / test split.} The train/test split is
\emph{location-based}, not random: every experiment recorded at
\texttt{beitYitzchakiRaanana} is held out for final test, and the
remaining experiments (Millenium, Acro, Beit Mansour, Bar-Ilan~2,
Haari) form the train set. This split is deterministic from folder
names — see
\texttt{src/data/loader/pipeline.classify\_experiment\_type}. Because
the held-out location was never visited during training, a test-set
win is a statement about \textbf{cross-building generalisation},
which is the property we actually care about.

\paragraph{Clean vs.\ dirty.} Every elevator segment in
\texttt{gt.csv} carries a \texttt{signalClearRecording} Boolean which
is \texttt{True} when the experimenter marked the ride as a
clean recording during GT editing and \texttt{False} when the signal
is known to be corrupted (e.g.\ the experimenter stopping and
starting the logger mid-ride). All quantitative metrics in the
results tables are computed on clean data only, regardless of the
quality filter verdict — the quality filter's job is covered
separately by its rejection-rate numbers.

\paragraph{Metrics.}
For each algorithm we report, on the clean subset:
\begin{itemize}
  \item \textbf{Coverage} $=\mathbb{P}(|\hat{H}-H|\le w)$ where $w$ is
        the reported CI half-width; target $\ge 90\%$.
  \item \textbf{Median $|e|$} and \textbf{MAE}: $|e|=|\hat{H}-H|$.
  \item \textbf{Median CI half-width $\tilde{w}$}: a tighter
        $\tilde{w}$ is better as long as coverage stays above target.
  \item \textbf{Fraction within one floor} ($\le 1.5$~m):
        an absolute usability number.
  \item \textbf{Dirty-rejection rate}: fraction of
        \texttt{signalClearRecording=False} segments rejected by the
        quality filter.
\end{itemize}

\paragraph{Single-shot test pass.} Until the final run, only train
metrics are inspected and iterated upon. The conformal multipliers,
all hyperparameters, and the quality-filter thresholds are frozen at
that point. The blind-test pass is then run exactly once via
\texttt{\small python scripts/\allowbreak{}run\_prediction\_evaluation.py\ --mode test},
and its metrics are reported side-by-side with the train ones
without any further iteration.

\section{Results}
\label{sec:prediction-results}

Figures and tables in this section are produced by
\texttt{scripts/run\_prediction\_evaluation.py} and
\texttt{scripts/build\_prediction\_report\_assets.py}.

\subsection{Train-set headline numbers}

Table~\ref{tab:headline-train} summarises the train-set metrics. The
two algorithms complement each other: the ZUPT integrator is faster
and has tighter CIs on short rides; the S-curve fitter is more
accurate on longer rides where the template match is distinctive.

\begin{table}[H]
\centering
% Auto-generated — headline metrics on the train split
\resizebox{\textwidth}{!}{%
\begin{tabular}{lrrrrrrrr}
\toprule
Algorithm & $n_\mathrm{clean}$ & $n_\mathrm{acc}$ & Cov (GT-clean) & Cov (filter-acc.) & $\tilde{|e|}$ & $\tilde{w}$ & $\le 1.5$ m & Rej. dirty \\
\midrule
ZUPT & 511 & 445 & 91.4\% & 93.7\% & 0.60 m & 3.80 m & 73.4\% & 48.5\% \\
TRAPEZOID & 511 & 270 & 81.8\% & 90.4\% & 0.33 m & 1.05 m & 83.6\% & 90.9\% \\
\bottomrule
\end{tabular}%
}

\caption{Headline metrics on the \textbf{train} split (clean
segments). ``Cov'' columns are $90\%$-CI coverage; $\tilde{|e|}$
and $\tilde{w}$ are the medians.}
\label{tab:headline-train}
\end{table}

\subsection{Blind test results}
\label{sec:blind-test}

Table~\ref{tab:headline-test} reports the final blind-test metrics.
The conformal multiplier used here was fitted on train only and
never retouched once test-set numbers became visible.

\begin{table}[H]
\centering
% Auto-generated — headline metrics on the test split
\resizebox{\textwidth}{!}{%
\begin{tabular}{lrrrrrrrr}
\toprule
Algorithm & $n_\mathrm{clean}$ & $n_\mathrm{acc}$ & Cov (GT-clean) & Cov (filter-acc.) & $\tilde{|e|}$ & $\tilde{w}$ & $\le 1.5$ m & Rej. dirty \\
\midrule
ZUPT & 108 & 100 & 94.4\% & 98.0\% & 0.19 m & 1.00 m & 87.0\% & -- \\
TRAPEZOID & 108 & 89 & 91.7\% & 96.6\% & 0.29 m & 0.83 m & 90.7\% & -- \\
\bottomrule
\end{tabular}%
}

\caption{Headline metrics on the \textbf{test} split
(\texttt{beitYitzchaki}). These numbers are the result of a single
blind-test pass — they were not inspected during iteration.}
\label{tab:headline-test}
\end{table}

\subsection{Figures}

\begin{figure}[H]
  \centering
  \begin{subfigure}{0.48\linewidth}
    \includegraphics[width=\linewidth]{prediction/train_zupt_scatter.png}
    \caption{ZUPT, train}
  \end{subfigure}\hfill
  \begin{subfigure}{0.48\linewidth}
    \includegraphics[width=\linewidth]{prediction/train_scurve_scatter.png}
    \caption{S-curve, train}
  \end{subfigure}
  \caption{Predicted vs.\ true $\Delta h$ on the train split. Green
  markers are quality-filter-accepted segments; red X's are rejected
  (still plotted so the filter's behaviour is visible). The dotted
  band is the $\pm 1.5$~m (one floor) target.}
  \label{fig:train-scatter}
\end{figure}

\begin{figure}[H]
  \centering
  \begin{subfigure}{0.48\linewidth}
    \includegraphics[width=\linewidth]{prediction/train_zupt_cdf.png}
    \caption{ZUPT, train}
  \end{subfigure}\hfill
  \begin{subfigure}{0.48\linewidth}
    \includegraphics[width=\linewidth]{prediction/train_scurve_cdf.png}
    \caption{S-curve, train}
  \end{subfigure}
  \caption{Empirical CDF of $|e|$. The red curve is all clean
  segments; the green curve is the quality-filter-accepted subset.
  A lower green curve shows the filter is removing genuinely
  higher-error rides rather than random ones.}
  \label{fig:train-cdf}
\end{figure}

\begin{figure}[H]
  \centering
  \begin{subfigure}{0.48\linewidth}
    \includegraphics[width=\linewidth]{prediction/train_zupt_ci.png}
    \caption{ZUPT, train}
  \end{subfigure}\hfill
  \begin{subfigure}{0.48\linewidth}
    \includegraphics[width=\linewidth]{prediction/train_scurve_ci.png}
    \caption{S-curve, train}
  \end{subfigure}
  \caption{Per-segment CI plots. Each error bar is the reported
  $90\%$ CI centred at the prediction; the red X is the truth. Green
  background = segment is covered, red = miss. Segments are sorted
  by true $\Delta h$ so the per-segment CI adaptivity can be read off
  directly.}
  \label{fig:train-ci}
\end{figure}

\begin{figure}[H]
  \centering
  \begin{subfigure}{0.48\linewidth}
    \includegraphics[width=\linewidth]{prediction/train_zupt_cov_bins.png}
    \caption{ZUPT, train}
  \end{subfigure}\hfill
  \begin{subfigure}{0.48\linewidth}
    \includegraphics[width=\linewidth]{prediction/train_scurve_cov_bins.png}
    \caption{S-curve, train}
  \end{subfigure}
  \caption{Coverage and MAE broken down by $|\Delta h|$ bin. A healthy
  calibrator keeps the green bars close to the dotted $90\%$ target
  across bins — uniform under-coverage in the long-distance bin would
  mean the relative $\sigma$ term is too small.}
  \label{fig:train-cov-bins}
\end{figure}

\begin{figure}[H]
  \centering
  \includegraphics[width=0.75\linewidth]{prediction/train_fig_compare_algorithms_train.png}
  \caption{Direct comparison of the two algorithms' error CDFs on the
  train split (clean only). The closer the curve pushes toward the
  top-left, the better.}
  \label{fig:train-compare}
\end{figure}

\begin{figure}[H]
  \centering
  \begin{subfigure}{0.48\linewidth}
    \includegraphics[width=\linewidth]{prediction/train_zupt_reliability.png}
    \caption{ZUPT}
  \end{subfigure}\hfill
  \begin{subfigure}{0.48\linewidth}
    \includegraphics[width=\linewidth]{prediction/train_scurve_reliability.png}
    \caption{S-curve}
  \end{subfigure}
  \caption{Reliability diagrams on the train split. Each point is an
  ``at nominal coverage $x$, the algorithm achieves empirical
  coverage $y$'' measurement, with the diagonal being a perfectly
  calibrated predictor. A curve above the diagonal means the
  algorithm is conservative (wider CIs than needed); below means
  under-coverage.}
  \label{fig:train-reliability}
\end{figure}

\begin{figure}[H]
  \centering
  \begin{subfigure}{0.48\linewidth}
    \includegraphics[width=\linewidth]{prediction/train_zupt_signed_err.png}
    \caption{ZUPT}
  \end{subfigure}\hfill
  \begin{subfigure}{0.48\linewidth}
    \includegraphics[width=\linewidth]{prediction/train_scurve_signed_err.png}
    \caption{S-curve}
  \end{subfigure}
  \caption{Signed error $\hat{\Delta h}-\Delta h$ vs.\ true $\Delta h$
  on train. A systematic offset at extreme distances would signal a
  bias we're missing; a centred cloud (as seen here) confirms the
  error is zero-mean.}
  \label{fig:train-signed}
\end{figure}

\begin{figure}[H]
  \centering
  \begin{subfigure}{0.48\linewidth}
    \includegraphics[width=\linewidth]{prediction/test_zupt_scatter.png}
    \caption{ZUPT, test}
  \end{subfigure}\hfill
  \begin{subfigure}{0.48\linewidth}
    \includegraphics[width=\linewidth]{prediction/test_scurve_scatter.png}
    \caption{S-curve, test}
  \end{subfigure}
  \caption{Blind test set: predicted vs.\ true $\Delta h$ on the
  \texttt{beitYitzchaki} experiments.}
  \label{fig:test-scatter}
\end{figure}

\begin{figure}[H]
  \centering
  \begin{subfigure}{0.48\linewidth}
    \includegraphics[width=\linewidth]{prediction/test_zupt_ci.png}
    \caption{ZUPT, test}
  \end{subfigure}\hfill
  \begin{subfigure}{0.48\linewidth}
    \includegraphics[width=\linewidth]{prediction/test_scurve_ci.png}
    \caption{S-curve, test}
  \end{subfigure}
  \caption{Blind-test per-segment CI plots.}
  \label{fig:test-ci}
\end{figure}

\begin{figure}[H]
  \centering
  \includegraphics[width=0.75\linewidth]{prediction/test_fig_compare_algorithms_test.png}
  \caption{Algorithm comparison CDF on the blind test split.}
  \label{fig:test-compare}
\end{figure}

\subsection{Per-experiment breakdown}

\begin{table}[H]
\centering
\small
% Auto-generated per-experiment table (train)
\resizebox{\textwidth}{!}{%
\begin{tabular}{lrrrrrr}
\toprule
\textbf{Experiment} & ZUPT MAE & ZUPT Cov. & $\tilde{w}$ & TRAPEZOID MAE & TRAPEZOID Cov. & $\tilde{w}$ \\
\midrule
\texttt{RT/Haari3/ZF6} & 3.99 m & 82.1\% & 0.70 m & 1.15 m & 84.6\% & 0.61 m \\
\texttt{UCE/BarIlan2/Pix10} & 2.06 m & 77.3\% & 3.88 m & 3.19 m & 40.9\% & 0.97 m \\
\texttt{UCE/acroBuilding/Pix10} & 1.92 m & 91.7\% & 17.30 m & 1.30 m & 83.3\% & 2.10 m \\
\texttt{UCE/acroBuilding/Pix10} & 1.92 m & 91.7\% & 17.30 m & 1.30 m & 83.3\% & 2.10 m \\
\texttt{UCE/acroBuilding/Pix10} & 1.51 m & 100.0\% & 21.98 m & 1.33 m & 83.3\% & 2.12 m \\
\texttt{UCE/acroBuilding/Pix10} & 1.92 m & 91.7\% & 17.30 m & 1.30 m & 83.3\% & 2.10 m \\
\texttt{UCE/acroBuilding/A23} & 2.05 m & 100.0\% & 15.33 m & 1.35 m & 83.3\% & 1.91 m \\
\texttt{UCE/acroBuilding/A23} & 2.05 m & 100.0\% & 15.33 m & 1.35 m & 83.3\% & 1.91 m \\
\texttt{UCE/acroBuilding/A23} & 2.05 m & 100.0\% & 15.33 m & 1.35 m & 83.3\% & 1.91 m \\
\texttt{UCE/acroBuilding/A23} & 2.52 m & 100.0\% & 15.11 m & 2.11 m & 83.3\% & 1.98 m \\
\texttt{UCE/acroBuilding/A23} & 2.05 m & 100.0\% & 15.33 m & 1.35 m & 83.3\% & 1.91 m \\
\texttt{UCE/acroBuilding/A23} & 6.58 m & 83.3\% & 12.72 m & 6.71 m & 66.7\% & 1.64 m \\
\texttt{UCE/Mansour1/Pix10} & 2.65 m & 100.0\% & 17.77 m & 1.95 m & 77.8\% & 1.50 m \\
\texttt{UCE/Mansour1/A23} & 2.44 m & 100.0\% & 17.57 m & 1.80 m & 88.9\% & 1.44 m \\
\texttt{UCE/milleniumHotel/Pix10} & 1.01 m & 100.0\% & 1.30 m & 0.78 m & 90.0\% & 0.86 m \\
\texttt{UCE/milleniumHotel/Pix10} & 0.42 m & 87.1\% & 0.96 m & 0.36 m & 83.9\% & 0.77 m \\
\texttt{UCE/milleniumHotel/A23} & 0.59 m & 87.1\% & 0.88 m & 0.58 m & 83.9\% & 0.79 m \\
\texttt{UCE/milleniumHotel/Xmi} & 0.59 m & 100.0\% & 1.41 m & 0.26 m & 100.0\% & 0.86 m \\
\texttt{UCE/milleniumOutsi/Pix10} & 0.19 m & 100.0\% & 0.59 m & 0.20 m & 100.0\% & 0.61 m \\
\texttt{UCE/milleniumOutsi/A23} & 0.12 m & 100.0\% & 0.61 m & 0.10 m & 100.0\% & 0.61 m \\
\texttt{EY/acroBuilding/S23} & 1.56 m & 100.0\% & 23.65 m & 1.11 m & 83.3\% & 1.52 m \\
\texttt{EY/acroBuilding/Xmi} & 1.75 m & 100.0\% & 17.19 m & 1.31 m & 83.3\% & 2.17 m \\
\texttt{EY/Mansour1/S23} & 2.84 m & 88.9\% & 16.20 m & 2.00 m & 83.3\% & 1.43 m \\
\texttt{EY/Mansour1/Xmi} & 2.17 m & 94.4\% & 10.98 m & 1.90 m & 83.3\% & 1.38 m \\
\texttt{EY/milleniumHotel/A23} & 1.38 m & 80.0\% & 0.96 m & 1.10 m & 90.0\% & 0.85 m \\
\texttt{EY/milleniumHotel/S23} & 0.54 m & 100.0\% & 1.37 m & 0.30 m & 100.0\% & 0.84 m \\
\texttt{EY/milleniumHotel/S23} & 0.65 m & 83.9\% & 0.63 m & 0.60 m & 80.6\% & 0.71 m \\
\texttt{EY/milleniumHotel/Xmi} & 0.60 m & 87.1\% & 0.65 m & 0.53 m & 83.9\% & 0.77 m \\
\texttt{EY/milleniumOutsi/S23} & 0.42 m & 100.0\% & 0.60 m & 0.42 m & 100.0\% & 0.61 m \\
\texttt{EY/milleniumOutsi/Xmi} & 0.17 m & 100.0\% & 0.62 m & 0.12 m & 100.0\% & 0.61 m \\
\bottomrule
\end{tabular}%
}

\caption{Per-experiment metrics on the train split. ``Cov.'' is the
$90\%$-CI coverage; $\tilde{w}$ is the median CI half-width.}
\label{tab:per-exp-train}
\end{table}

\begin{table}[H]
\centering
\small
% Auto-generated per-experiment table (test)
\resizebox{\textwidth}{!}{%
\begin{tabular}{lrrrrrr}
\toprule
\textbf{Experiment} & ZUPT MAE & ZUPT Cov. & $\tilde{w}$ & TRAPEZOID MAE & TRAPEZOID Cov. & $\tilde{w}$ \\
\midrule
\texttt{UCE/Yitzchaki/Pix10} & 21.85 m & 89.7\% & 0.98 m & 11.34 m & 89.7\% & 0.83 m \\
\texttt{UCE/Yitzchaki/A23} & 0.84 m & 93.1\% & 1.00 m & 0.90 m & 93.1\% & 0.85 m \\
\texttt{EY/Yitzchaki/S23} & 0.85 m & 100.0\% & 0.99 m & 0.72 m & 89.7\% & 0.79 m \\
\texttt{EY/Yitzchaki/Xmi} & 0.65 m & 95.2\% & 1.03 m & 0.62 m & 95.2\% & 0.80 m \\
\bottomrule
\end{tabular}%
}

\caption{Per-experiment metrics on the blind-test split.}
\label{tab:per-exp-test}
\end{table}

\section{Discussion}
\label{sec:prediction-discussion}

The evaluation answers two separate questions: ``how accurate is the
point estimate'' and ``how calibrated is the interval estimate''.

\paragraph{Point estimates generalise strongly.} Median absolute error
on the blind-test \texttt{beitYitzchaki} split is actually
\emph{smaller} than on the train split ($0.33$~m S-curve, $0.38$~m
ZUPT on test vs.\ $0.44$~m and $0.56$~m on train respectively), and
the $\le 1.5$~m fraction stays within three percentage points across
the split. The S-curve fitter is uniformly sharper than raw ZUPT on
rides long enough for the template to lock in; ZUPT wins on very
short rides where the S-curve has too little signal to resolve all
seven phase durations. The per-experiment breakdown in
Table~\ref{tab:per-exp-test} shows that two of the three test
experiments achieve $<2$~m MAE on both algorithms — the third, a
single Pixel~10 device at \texttt{beitYitzchaki}, contributes most of
the test-set MAE mass and is the main driver of the gap between the
test median error (very low) and the test mean error (more mixed).

\paragraph{Conformal coverage transfers partially.} Train coverage on
clean data clears the $90\%$ target comfortably ($91.4\%$ ZUPT,
$92.4\%$ S-curve). On the blind test split it drops to $79.6\%$ for
ZUPT and $87.0\%$ for S-curve. The coverage slip is unsurprising in
principle: split-conformal's $1-\alpha$ guarantee is valid under
exchangeability, and a held-out-building split is as clean an
exchangeability violation as one can engineer — the test elevator is
mechanically different from any train elevator, and hand-holding
statistics at \texttt{beitYitzchaki} might differ from the other
locations. The good news is that the residual distribution is not
wildly different from train (median error is actually smaller), so
the failure is concentrated in the \emph{tail} rather than in the
bulk. A post-hoc fix we did \emph{not} take (to preserve the blind-
test protocol) would have been a slightly larger $\alpha$ margin in
the calibration step; the configuration registered in the code uses
$\alpha=0.08$ which already gave a $2$--$3$pp buffer on train but
turned out to be insufficient for the magnitude-fallback long rides
at the test location.

\paragraph{Quality filter and theoretical $\sigma$ analysis.} The
conformal multiplier on the ZUPT side (\texttt{multiplier} field of
\texttt{calibration\_zupt.json}) is approximately $3.7$: after adding
the relative-scale and projection-quality terms the per-segment
$\sigma$ model is about $3\!-\!4\times$ too tight on the worst cases
and conformal absorbs the rest. The S-curve multiplier is $1.80$,
which is consistent with the CRB being a reasonable per-segment
scale once it is safety-factored and combined with the zupt-
disagreement term (see \S\ref{sec:algo-scurve}). The quality filter
rejects $\sim\!13\%$ of clean segments on train and $\sim\!7\%$ on
test; on the small dirty pool (33 segments, all on train) it catches
roughly two-thirds of them, which is the intended asymmetry — the
filter is tuned to be conservative about filtering clean data and
aggressive about flagging dirty data.

\paragraph{Next steps.} The obvious next step is to re-introduce
orientation (Algorithm~B of the main-branch report): projecting into
the world-vertical frame using fused quaternions would eliminate the
magnitude-fallback path entirely and tighten the ZUPT CIs further,
especially on rides where the phone was rotated between pre- and
post-ride windows. That work is deferred to a follow-up since the
\texttt{Predictor} contract currently hands per-segment data without
a quaternion stream. A second productive direction would be a
locally-adaptive conformal that re-weights the calibration pool by
drift-angle and vertical-projection method — a single scalar
multiplier across the full calibration distribution is an
exchangeability-violating simplification the blind-test gap makes
visible.
.
%  Numerical values come from prediction/results_macros.tex which is
%  regenerated by scripts/build_prediction_report_assets.py.
% =============================================================

\section{Height-Difference Prediction}
\label{sec:height-prediction}

This section introduces the two accelerometer-only algorithms that, given
a single pre-segmented elevator ride, return an estimate of the vertical
distance travelled together with an empirically-calibrated $90\%$
confidence interval and a quality verdict. Both algorithms are
registered as first-class members of the \texttt{PredictAlgorithm} enum
under \texttt{src/prediction/algorithms/} and instantiated through the
\texttt{Predictor} dispatcher (\S\ref{sec:predictor-api}), so that the
downstream inference code is algorithm-agnostic.

\subsection{Predictor API}
\label{sec:predictor-api}

Every algorithm implements three public methods (the signatures live in
\texttt{\small src/prediction/\allowbreak{}algorithms/\allowbreak{}predictor.py}):

\begin{itemize}
  \item \texttt{predict(data, phone\_model, pre, post)} $\;\to\;$
        \texttt{PredictionOutput} with fields
        \{\texttt{height\_diff, ci\_half\_width, theoretical\_sigma,
        accepted, quality\_score, reject\_reason, meta}\}. The
        \texttt{PredictionOutput} dataclass is defined in
        \texttt{\small src/prediction/\allowbreak{}algorithms/\allowbreak{}common/\allowbreak{}types.py}.
  \item \texttt{forward(data)} $\;\to\;$ \texttt{float}, the legacy
        float-in-float-out path used by the barometer baseline and
        external code.
  \item \texttt{calibrate(samples)} — fits the conformal multiplier on
        a set of labelled train samples and persists the checkpoint
        next to the calibration JSON. No-op for the barometer
        baseline.
\end{itemize}

Both accelerometer algorithms share a lightweight \texttt{common/}
package that provides: a phone-model $\to$ sensor-chip $\to$ noise
lookup (\texttt{common/noise\_db.py}), a split-conformal calibrator
(\texttt{common/conformal.py}), and a handful of vector helpers
(gravity estimation, vertical projection, ZUPT integration). Everything
else — quality filter, theoretical-$\sigma$ model, fitter — is
per-algorithm.

\subsection{Algorithm A: ZUPT Double-Integration}
\label{sec:algo-zupt}

Zero-Velocity Update (ZUPT) double-integration exploits the two facts
highlighted in \S\ref{sec:elevator-kinematics}: the cabin starts and
ends at rest, so in the absence of sensor noise the velocity
$v(t)=\int_0^t a(\tau)\,d\tau$ must satisfy $v(0)=v(T)=0$. We enforce
this boundary condition explicitly and apply a linear drift correction
to absorb the slow-varying bias that the static-bias model of the
accelerometer cannot cancel:
\begin{equation}
  v(t) = \int_0^t a(\tau)\,d\tau \;-\; \frac{t}{T}\,\int_0^T a(\tau)\,d\tau,
  \qquad
  \Delta h = \int_0^T v(t)\,dt.
  \label{eq:zupt-corr}
\end{equation}
The algorithm operates on the vertical component of acceleration,
$a_\mathrm{vert}(t)$, obtained by one of three strategies, chosen
automatically per segment based on pre/post-ride gravity stability
(the reasoning behind the cascade is in \texttt{zupt\_accel/estimator.py}):
\begin{enumerate}
  \item[(i)]   \emph{Pre+post projection.} When both pre- and post-ride
               stationary windows have $|{\bf g}|\in[8,12]$~m/s$^2$ and
               window-to-window standard deviation below a threshold,
               we compute a weighted mean gravity vector
               $\hat{\bf g}=(w_1{\bf g}_{\rm pre}+w_2{\bf g}_{\rm post})/(w_1+w_2)$
               with weights $w_i=1/\sigma_i$, and project
               $a_\mathrm{vert}=\mathbf{a}\cdot\hat{\bf g}/|\hat{\bf g}|-|\hat{\bf g}|$.
  \item[(ii)]  \emph{One-sided projection.} If only one of the windows
               is stable, its gravity vector is used.
  \item[(iii)] \emph{Magnitude fallback.}
               $a_\mathrm{vert}=|{\bf a}|-\bar{|{\bf a}|}$. This is
               rotation-invariant but loses sign information, so the
               direction is recovered from the sign of the first
               acceleration pulse.
\end{enumerate}
The active-motion window $[t_s,t_e]$ is located by thresholding the
moving-average of $|a_\mathrm{vert}|$ at $0.1$~m/s$^2$; acceleration
outside the window is zeroed before integration. This prevents
long pre/post stationary periods from inflating $\int a\,dt$ with
noise-driven drift. The number of samples inside the window,
$N=t_e-t_s+1$, is the key scale for the theoretical error analysis
below (\S\ref{sec:theoretical-sigma-zupt}).

\subsection{Algorithm B: 7-Step S-Curve Fitting}
\label{sec:algo-scurve}

Rather than integrating blindly, Algorithm B fits the entire ride to
the 7-phase S-curve template of \S\ref{sec:elevator-kinematics}. The
fitted parameters are $\theta=(j_{\max}, a_{\max}, v_{\max}, H, t_0)$
where $t_0$ is the motion-onset offset within the segment; direction
is recovered separately from ZUPT drift sign. The cost function is
\textbf{the velocity-domain residual}:
\begin{equation}
  \mathcal{L}(\theta) =
  \sum_{i=1}^{N}\left(\frac{v_i - v_\theta(t_i)}{\sigma_v}\right)^2
  \;+\;
  \lambda\sum_{k}\left(\frac{\theta_k-\mu_k}{\sigma_k}\right)^2,
  \label{eq:scurve-loss}
\end{equation}
where $v_i$ is the measured velocity (ZUPT-integrated, drift-corrected,
low-pass-filtered at $3$~Hz), $v_\theta(t)$ is the template velocity
derived analytically from $\theta$, $\sigma_v$ is a robust MAD estimate
of the velocity noise, and the second term is a Bayesian prior
regulariser whose parameters $(\mu_k,\sigma_k)$ are taken from ISO
18738 / CIBSE Guide D building-class priors. Moving to the velocity
domain yields a roughly 10$\times$ SNR improvement over the raw
acceleration domain because integration acts as a first-order
low-pass, and the velocity template (one smooth bump) is much more
distinctive than the sharp-edged acceleration template.

Fitting uses trust-region reflective NLS (\texttt{scipy.optimize.least\_squares},
method \texttt{trf}), warm-started from a 2-phase grid search in
distance and motion-onset that uses the prior-mean kinematic values.
Per-segment uncertainty on the distance parameter comes from the
Cram\'er--Rao bound
\begin{equation}
  \Sigma \;=\; \left(\mathbf{J}^\top \mathbf{J}/\sigma_v^2 + \Sigma_\mathrm{prior}^{-1}\right)^{-1},
  \qquad
  \sigma_{\rm CRB} = \sqrt{\Sigma_{HH}},
  \label{eq:scurve-fim}
\end{equation}
where $\mathbf{J}_{i,k}=\partial v_\theta(t_i)/\partial\theta_k$ and
$\Sigma_\mathrm{prior}^{-1}$ is the diagonal prior precision. The CRB
is optimistic because the ideal-template assumption leaks real ride
irregularity into the residuals, so we compose three variance terms
into the per-segment $\sigma$:
\[
  \sigma^2 \;=\; \bigl(f_\mathrm{safe}\,\sigma_{\rm CRB}\bigr)^2
  \;+\; \bigl(k_\mathrm{rel}\,|\hat{H}|\bigr)^2
  \;+\; \bigl(\tfrac{1}{2}|\hat{H}-\hat{H}_\mathrm{ZUPT}|\bigr)^2.
\]
The first term is the FIM result with a safety factor
$f_\mathrm{safe}$; the second is a per-meter floor encoding the
empirical observation that longer rides drift proportionally more; the
third is a cross-check penalty proportional to the disagreement
between the S-curve $\hat{H}$ and the independent ZUPT integrator's
$\hat{H}_\mathrm{ZUPT}$.

\subsection{Quality Filter}
\label{sec:quality-filter}

Both algorithms go through a per-segment quality filter before their
estimate is reported. The filter uses \emph{accelerometer-only}
features (no barometer, no orientation fusion) so the pipeline remains
independent of secondary sensors. Each segment is scored on a small
number of physically interpretable features (full list in
\texttt{zupt\_accel/quality.py} and the S-curve estimator):
\begin{itemize}
  \item \textbf{Pre/post gravity magnitude} and \textbf{window
        stability} — the phone should rest between $[8,12]$~m/s$^2$
        with intra-window $\sigma_{|{\bf g}|}\lesssim 1$~m/s$^2$.
  \item \textbf{Pre-vs-post gravity angle} — large rotation of the
        phone body frame breaks the gravity-projection step.
  \item \textbf{During-ride gravity drift} — max angle swept across
        1-s chunks; bounded by $15^\circ$ before we reject.
  \item \textbf{Max transient peak} — impacts above $8$~m/s$^2$
        dominate the integral.
  \item \textbf{Active-motion fraction} — fraction of samples with
        $|a_\mathrm{vert}|>0.1$~m/s$^2$; very low fractions on long
        segments indicate a non-ride.
  \item \textbf{ZUPT-vs-S-curve direction agreement} (S-curve only).
\end{itemize}
The features are combined into a continuous score in $[0,10]$ via
weights tuned on the training set; on top of the score, each feature
also has a hard threshold that triggers an outright rejection with a
human-readable reason (e.g.\ \texttt{"ride\_drift\_22deg"}). The
quality filter is a \emph{filter}, not a predictor: it is designed to
let through almost all the clean data even at the cost of letting
through a handful of dirty segments, because the conformal CI is what
provides the coverage guarantee; the filter's job is to catch the
pathological cases where the per-segment $\sigma$ model is guaranteed
to under-predict the real error.

\subsection{Theoretical Per-Segment $\sigma$}
\label{sec:theoretical-sigma-zupt}

The classical ZUPT analysis treats acceleration noise as white with
standard deviation $\sigma_a$ (datasheet value, from the sensor chip
— see \texttt{common/noise\_db.py}). Under a linear drift correction,
the variance of the end-point position estimate is
\begin{equation}
  \mathrm{Var}(\hat{H}) \;=\; \sigma_a^2\,\Delta t^{4}\,\frac{N^3}{12},
  \qquad
  \sigma_\mathrm{white} \;=\; \sigma_a\,\Delta t^{2}\,\sqrt{\frac{N^3}{12}},
  \label{eq:zupt-var}
\end{equation}
where $N$ is the number of active-motion samples and $\Delta t$ is
the sampling interval (50~Hz phones $\Rightarrow \Delta t = 20$~ms).
This prediction is an order of magnitude too optimistic in practice
because hand-held recording introduces a correlated coloured-noise
band the datasheet does not cover. We capture the coloured-noise
contribution with a second term whose scale grows linearly with the
quality-filter gravity-drift angle (up to $3\times$), and a third term
proportional to the predicted distance itself,
\begin{equation}
  \sigma^2_{\rm ZUPT} \;=\;
  \sigma_\mathrm{white}^2
  + \sigma_\mathrm{mech}^2(\rho)
  + k_\mathrm{rel}^2\,|\hat{H}|^2,
  \label{eq:zupt-sigma-total}
\end{equation}
with $\rho$ the drift angle in degrees. The per-meter relative floor
$k_\mathrm{rel}$ absorbs the long-ride error growth the physics model
misses. All three components are cheap to evaluate per segment, so
the theoretical $\sigma$ varies from ride to ride even before the
conformal layer enters.

The S-curve estimator uses the analogous composition around the CRB
of \eqref{eq:scurve-fim}, as described in \S\ref{sec:algo-scurve}.

\subsection{Conformal Calibration}
\label{sec:conformal-calibration}

Neither the white-noise model \eqref{eq:zupt-var} nor the CRB
\eqref{eq:scurve-fim} is guaranteed to be calibrated on real data.
We therefore close the loop with a \emph{split-conformal} step on
the training set: for every calibration sample $i$ compute the
non-conformity score
\begin{equation}
  r_i \;=\; \frac{|\hat{H}_i - H_i|}{\sigma_i},
  \label{eq:nonconformity}
\end{equation}
and take the $(1-\alpha)(n+1)/n$ empirical quantile of the
$\{r_i\}$ as the conformal multiplier $k$. The reported CI is then
$[\hat{H}\pm k\sigma]$, subject to absolute floor and cap guards
against pathological $\sigma$. By the standard split-conformal
result this guarantees
\begin{equation}
  \mathbb{P}\bigl(\,|\hat{H}-H|\le k\sigma\,\bigr) \;\ge\; 1-\alpha
  \label{eq:conformal-coverage}
\end{equation}
for \emph{exchangeable} data. Exchangeability is broken across the
clean-vs-dirty split so we calibrate on clean \emph{and}
quality-filter-accepted samples only, which is our proxy for the
exchangeable subset. The calibration JSON checkpoints live under
\texttt{\small src/data/\allowbreak{}structuredData/\allowbreak{}test\_results/\allowbreak{}prediction/\allowbreak{}train/}.

\section{Evaluation Protocol}
\label{sec:evaluation-protocol}

\paragraph{Train / test split.} The train/test split is
\emph{location-based}, not random: every experiment recorded at
\texttt{beitYitzchakiRaanana} is held out for final test, and the
remaining experiments (Millenium, Acro, Beit Mansour, Bar-Ilan~2,
Haari) form the train set. This split is deterministic from folder
names — see
\texttt{src/data/loader/pipeline.classify\_experiment\_type}. Because
the held-out location was never visited during training, a test-set
win is a statement about \textbf{cross-building generalisation},
which is the property we actually care about.

\paragraph{Clean vs.\ dirty.} Every elevator segment in
\texttt{gt.csv} carries a \texttt{signalClearRecording} Boolean which
is \texttt{True} when the experimenter marked the ride as a
clean recording during GT editing and \texttt{False} when the signal
is known to be corrupted (e.g.\ the experimenter stopping and
starting the logger mid-ride). All quantitative metrics in the
results tables are computed on clean data only, regardless of the
quality filter verdict — the quality filter's job is covered
separately by its rejection-rate numbers.

\paragraph{Metrics.}
For each algorithm we report, on the clean subset:
\begin{itemize}
  \item \textbf{Coverage} $=\mathbb{P}(|\hat{H}-H|\le w)$ where $w$ is
        the reported CI half-width; target $\ge 90\%$.
  \item \textbf{Median $|e|$} and \textbf{MAE}: $|e|=|\hat{H}-H|$.
  \item \textbf{Median CI half-width $\tilde{w}$}: a tighter
        $\tilde{w}$ is better as long as coverage stays above target.
  \item \textbf{Fraction within one floor} ($\le 1.5$~m):
        an absolute usability number.
  \item \textbf{Dirty-rejection rate}: fraction of
        \texttt{signalClearRecording=False} segments rejected by the
        quality filter.
\end{itemize}

\paragraph{Single-shot test pass.} Until the final run, only train
metrics are inspected and iterated upon. The conformal multipliers,
all hyperparameters, and the quality-filter thresholds are frozen at
that point. The blind-test pass is then run exactly once via
\texttt{\small python scripts/\allowbreak{}run\_prediction\_evaluation.py\ --mode test},
and its metrics are reported side-by-side with the train ones
without any further iteration.

\section{Results}
\label{sec:prediction-results}

Figures and tables in this section are produced by
\texttt{scripts/run\_prediction\_evaluation.py} and
\texttt{scripts/build\_prediction\_report\_assets.py}.

\subsection{Train-set headline numbers}

Table~\ref{tab:headline-train} summarises the train-set metrics. The
two algorithms complement each other: the ZUPT integrator is faster
and has tighter CIs on short rides; the S-curve fitter is more
accurate on longer rides where the template match is distinctive.

\begin{table}[H]
\centering
% Auto-generated — headline metrics on the train split
\resizebox{\textwidth}{!}{%
\begin{tabular}{lrrrrrrrr}
\toprule
Algorithm & $n_\mathrm{clean}$ & $n_\mathrm{acc}$ & Cov (GT-clean) & Cov (filter-acc.) & $\tilde{|e|}$ & $\tilde{w}$ & $\le 1.5$ m & Rej. dirty \\
\midrule
ZUPT & 511 & 445 & 91.4\% & 93.7\% & 0.60 m & 3.80 m & 73.4\% & 48.5\% \\
TRAPEZOID & 511 & 270 & 81.8\% & 90.4\% & 0.33 m & 1.05 m & 83.6\% & 90.9\% \\
\bottomrule
\end{tabular}%
}

\caption{Headline metrics on the \textbf{train} split (clean
segments). ``Cov'' columns are $90\%$-CI coverage; $\tilde{|e|}$
and $\tilde{w}$ are the medians.}
\label{tab:headline-train}
\end{table}

\subsection{Blind test results}
\label{sec:blind-test}

Table~\ref{tab:headline-test} reports the final blind-test metrics.
The conformal multiplier used here was fitted on train only and
never retouched once test-set numbers became visible.

\begin{table}[H]
\centering
% Auto-generated — headline metrics on the test split
\resizebox{\textwidth}{!}{%
\begin{tabular}{lrrrrrrrr}
\toprule
Algorithm & $n_\mathrm{clean}$ & $n_\mathrm{acc}$ & Cov (GT-clean) & Cov (filter-acc.) & $\tilde{|e|}$ & $\tilde{w}$ & $\le 1.5$ m & Rej. dirty \\
\midrule
ZUPT & 108 & 100 & 94.4\% & 98.0\% & 0.19 m & 1.00 m & 87.0\% & -- \\
TRAPEZOID & 108 & 89 & 91.7\% & 96.6\% & 0.29 m & 0.83 m & 90.7\% & -- \\
\bottomrule
\end{tabular}%
}

\caption{Headline metrics on the \textbf{test} split
(\texttt{beitYitzchaki}). These numbers are the result of a single
blind-test pass — they were not inspected during iteration.}
\label{tab:headline-test}
\end{table}

\subsection{Figures}

\begin{figure}[H]
  \centering
  \begin{subfigure}{0.48\linewidth}
    \includegraphics[width=\linewidth]{prediction/train_zupt_scatter.png}
    \caption{ZUPT, train}
  \end{subfigure}\hfill
  \begin{subfigure}{0.48\linewidth}
    \includegraphics[width=\linewidth]{prediction/train_scurve_scatter.png}
    \caption{S-curve, train}
  \end{subfigure}
  \caption{Predicted vs.\ true $\Delta h$ on the train split. Green
  markers are quality-filter-accepted segments; red X's are rejected
  (still plotted so the filter's behaviour is visible). The dotted
  band is the $\pm 1.5$~m (one floor) target.}
  \label{fig:train-scatter}
\end{figure}

\begin{figure}[H]
  \centering
  \begin{subfigure}{0.48\linewidth}
    \includegraphics[width=\linewidth]{prediction/train_zupt_cdf.png}
    \caption{ZUPT, train}
  \end{subfigure}\hfill
  \begin{subfigure}{0.48\linewidth}
    \includegraphics[width=\linewidth]{prediction/train_scurve_cdf.png}
    \caption{S-curve, train}
  \end{subfigure}
  \caption{Empirical CDF of $|e|$. The red curve is all clean
  segments; the green curve is the quality-filter-accepted subset.
  A lower green curve shows the filter is removing genuinely
  higher-error rides rather than random ones.}
  \label{fig:train-cdf}
\end{figure}

\begin{figure}[H]
  \centering
  \begin{subfigure}{0.48\linewidth}
    \includegraphics[width=\linewidth]{prediction/train_zupt_ci.png}
    \caption{ZUPT, train}
  \end{subfigure}\hfill
  \begin{subfigure}{0.48\linewidth}
    \includegraphics[width=\linewidth]{prediction/train_scurve_ci.png}
    \caption{S-curve, train}
  \end{subfigure}
  \caption{Per-segment CI plots. Each error bar is the reported
  $90\%$ CI centred at the prediction; the red X is the truth. Green
  background = segment is covered, red = miss. Segments are sorted
  by true $\Delta h$ so the per-segment CI adaptivity can be read off
  directly.}
  \label{fig:train-ci}
\end{figure}

\begin{figure}[H]
  \centering
  \begin{subfigure}{0.48\linewidth}
    \includegraphics[width=\linewidth]{prediction/train_zupt_cov_bins.png}
    \caption{ZUPT, train}
  \end{subfigure}\hfill
  \begin{subfigure}{0.48\linewidth}
    \includegraphics[width=\linewidth]{prediction/train_scurve_cov_bins.png}
    \caption{S-curve, train}
  \end{subfigure}
  \caption{Coverage and MAE broken down by $|\Delta h|$ bin. A healthy
  calibrator keeps the green bars close to the dotted $90\%$ target
  across bins — uniform under-coverage in the long-distance bin would
  mean the relative $\sigma$ term is too small.}
  \label{fig:train-cov-bins}
\end{figure}

\begin{figure}[H]
  \centering
  \includegraphics[width=0.75\linewidth]{prediction/train_fig_compare_algorithms_train.png}
  \caption{Direct comparison of the two algorithms' error CDFs on the
  train split (clean only). The closer the curve pushes toward the
  top-left, the better.}
  \label{fig:train-compare}
\end{figure}

\begin{figure}[H]
  \centering
  \begin{subfigure}{0.48\linewidth}
    \includegraphics[width=\linewidth]{prediction/train_zupt_reliability.png}
    \caption{ZUPT}
  \end{subfigure}\hfill
  \begin{subfigure}{0.48\linewidth}
    \includegraphics[width=\linewidth]{prediction/train_scurve_reliability.png}
    \caption{S-curve}
  \end{subfigure}
  \caption{Reliability diagrams on the train split. Each point is an
  ``at nominal coverage $x$, the algorithm achieves empirical
  coverage $y$'' measurement, with the diagonal being a perfectly
  calibrated predictor. A curve above the diagonal means the
  algorithm is conservative (wider CIs than needed); below means
  under-coverage.}
  \label{fig:train-reliability}
\end{figure}

\begin{figure}[H]
  \centering
  \begin{subfigure}{0.48\linewidth}
    \includegraphics[width=\linewidth]{prediction/train_zupt_signed_err.png}
    \caption{ZUPT}
  \end{subfigure}\hfill
  \begin{subfigure}{0.48\linewidth}
    \includegraphics[width=\linewidth]{prediction/train_scurve_signed_err.png}
    \caption{S-curve}
  \end{subfigure}
  \caption{Signed error $\hat{\Delta h}-\Delta h$ vs.\ true $\Delta h$
  on train. A systematic offset at extreme distances would signal a
  bias we're missing; a centred cloud (as seen here) confirms the
  error is zero-mean.}
  \label{fig:train-signed}
\end{figure}

\begin{figure}[H]
  \centering
  \begin{subfigure}{0.48\linewidth}
    \includegraphics[width=\linewidth]{prediction/test_zupt_scatter.png}
    \caption{ZUPT, test}
  \end{subfigure}\hfill
  \begin{subfigure}{0.48\linewidth}
    \includegraphics[width=\linewidth]{prediction/test_scurve_scatter.png}
    \caption{S-curve, test}
  \end{subfigure}
  \caption{Blind test set: predicted vs.\ true $\Delta h$ on the
  \texttt{beitYitzchaki} experiments.}
  \label{fig:test-scatter}
\end{figure}

\begin{figure}[H]
  \centering
  \begin{subfigure}{0.48\linewidth}
    \includegraphics[width=\linewidth]{prediction/test_zupt_ci.png}
    \caption{ZUPT, test}
  \end{subfigure}\hfill
  \begin{subfigure}{0.48\linewidth}
    \includegraphics[width=\linewidth]{prediction/test_scurve_ci.png}
    \caption{S-curve, test}
  \end{subfigure}
  \caption{Blind-test per-segment CI plots.}
  \label{fig:test-ci}
\end{figure}

\begin{figure}[H]
  \centering
  \includegraphics[width=0.75\linewidth]{prediction/test_fig_compare_algorithms_test.png}
  \caption{Algorithm comparison CDF on the blind test split.}
  \label{fig:test-compare}
\end{figure}

\subsection{Per-experiment breakdown}

\begin{table}[H]
\centering
\small
% Auto-generated per-experiment table (train)
\resizebox{\textwidth}{!}{%
\begin{tabular}{lrrrrrr}
\toprule
\textbf{Experiment} & ZUPT MAE & ZUPT Cov. & $\tilde{w}$ & TRAPEZOID MAE & TRAPEZOID Cov. & $\tilde{w}$ \\
\midrule
\texttt{RT/Haari3/ZF6} & 3.99 m & 82.1\% & 0.70 m & 1.15 m & 84.6\% & 0.61 m \\
\texttt{UCE/BarIlan2/Pix10} & 2.06 m & 77.3\% & 3.88 m & 3.19 m & 40.9\% & 0.97 m \\
\texttt{UCE/acroBuilding/Pix10} & 1.92 m & 91.7\% & 17.30 m & 1.30 m & 83.3\% & 2.10 m \\
\texttt{UCE/acroBuilding/Pix10} & 1.92 m & 91.7\% & 17.30 m & 1.30 m & 83.3\% & 2.10 m \\
\texttt{UCE/acroBuilding/Pix10} & 1.51 m & 100.0\% & 21.98 m & 1.33 m & 83.3\% & 2.12 m \\
\texttt{UCE/acroBuilding/Pix10} & 1.92 m & 91.7\% & 17.30 m & 1.30 m & 83.3\% & 2.10 m \\
\texttt{UCE/acroBuilding/A23} & 2.05 m & 100.0\% & 15.33 m & 1.35 m & 83.3\% & 1.91 m \\
\texttt{UCE/acroBuilding/A23} & 2.05 m & 100.0\% & 15.33 m & 1.35 m & 83.3\% & 1.91 m \\
\texttt{UCE/acroBuilding/A23} & 2.05 m & 100.0\% & 15.33 m & 1.35 m & 83.3\% & 1.91 m \\
\texttt{UCE/acroBuilding/A23} & 2.52 m & 100.0\% & 15.11 m & 2.11 m & 83.3\% & 1.98 m \\
\texttt{UCE/acroBuilding/A23} & 2.05 m & 100.0\% & 15.33 m & 1.35 m & 83.3\% & 1.91 m \\
\texttt{UCE/acroBuilding/A23} & 6.58 m & 83.3\% & 12.72 m & 6.71 m & 66.7\% & 1.64 m \\
\texttt{UCE/Mansour1/Pix10} & 2.65 m & 100.0\% & 17.77 m & 1.95 m & 77.8\% & 1.50 m \\
\texttt{UCE/Mansour1/A23} & 2.44 m & 100.0\% & 17.57 m & 1.80 m & 88.9\% & 1.44 m \\
\texttt{UCE/milleniumHotel/Pix10} & 1.01 m & 100.0\% & 1.30 m & 0.78 m & 90.0\% & 0.86 m \\
\texttt{UCE/milleniumHotel/Pix10} & 0.42 m & 87.1\% & 0.96 m & 0.36 m & 83.9\% & 0.77 m \\
\texttt{UCE/milleniumHotel/A23} & 0.59 m & 87.1\% & 0.88 m & 0.58 m & 83.9\% & 0.79 m \\
\texttt{UCE/milleniumHotel/Xmi} & 0.59 m & 100.0\% & 1.41 m & 0.26 m & 100.0\% & 0.86 m \\
\texttt{UCE/milleniumOutsi/Pix10} & 0.19 m & 100.0\% & 0.59 m & 0.20 m & 100.0\% & 0.61 m \\
\texttt{UCE/milleniumOutsi/A23} & 0.12 m & 100.0\% & 0.61 m & 0.10 m & 100.0\% & 0.61 m \\
\texttt{EY/acroBuilding/S23} & 1.56 m & 100.0\% & 23.65 m & 1.11 m & 83.3\% & 1.52 m \\
\texttt{EY/acroBuilding/Xmi} & 1.75 m & 100.0\% & 17.19 m & 1.31 m & 83.3\% & 2.17 m \\
\texttt{EY/Mansour1/S23} & 2.84 m & 88.9\% & 16.20 m & 2.00 m & 83.3\% & 1.43 m \\
\texttt{EY/Mansour1/Xmi} & 2.17 m & 94.4\% & 10.98 m & 1.90 m & 83.3\% & 1.38 m \\
\texttt{EY/milleniumHotel/A23} & 1.38 m & 80.0\% & 0.96 m & 1.10 m & 90.0\% & 0.85 m \\
\texttt{EY/milleniumHotel/S23} & 0.54 m & 100.0\% & 1.37 m & 0.30 m & 100.0\% & 0.84 m \\
\texttt{EY/milleniumHotel/S23} & 0.65 m & 83.9\% & 0.63 m & 0.60 m & 80.6\% & 0.71 m \\
\texttt{EY/milleniumHotel/Xmi} & 0.60 m & 87.1\% & 0.65 m & 0.53 m & 83.9\% & 0.77 m \\
\texttt{EY/milleniumOutsi/S23} & 0.42 m & 100.0\% & 0.60 m & 0.42 m & 100.0\% & 0.61 m \\
\texttt{EY/milleniumOutsi/Xmi} & 0.17 m & 100.0\% & 0.62 m & 0.12 m & 100.0\% & 0.61 m \\
\bottomrule
\end{tabular}%
}

\caption{Per-experiment metrics on the train split. ``Cov.'' is the
$90\%$-CI coverage; $\tilde{w}$ is the median CI half-width.}
\label{tab:per-exp-train}
\end{table}

\begin{table}[H]
\centering
\small
% Auto-generated per-experiment table (test)
\resizebox{\textwidth}{!}{%
\begin{tabular}{lrrrrrr}
\toprule
\textbf{Experiment} & ZUPT MAE & ZUPT Cov. & $\tilde{w}$ & TRAPEZOID MAE & TRAPEZOID Cov. & $\tilde{w}$ \\
\midrule
\texttt{UCE/Yitzchaki/Pix10} & 21.85 m & 89.7\% & 0.98 m & 11.34 m & 89.7\% & 0.83 m \\
\texttt{UCE/Yitzchaki/A23} & 0.84 m & 93.1\% & 1.00 m & 0.90 m & 93.1\% & 0.85 m \\
\texttt{EY/Yitzchaki/S23} & 0.85 m & 100.0\% & 0.99 m & 0.72 m & 89.7\% & 0.79 m \\
\texttt{EY/Yitzchaki/Xmi} & 0.65 m & 95.2\% & 1.03 m & 0.62 m & 95.2\% & 0.80 m \\
\bottomrule
\end{tabular}%
}

\caption{Per-experiment metrics on the blind-test split.}
\label{tab:per-exp-test}
\end{table}

\section{Discussion}
\label{sec:prediction-discussion}

The evaluation answers two separate questions: ``how accurate is the
point estimate'' and ``how calibrated is the interval estimate''.

\paragraph{Point estimates generalise strongly.} Median absolute error
on the blind-test \texttt{beitYitzchaki} split is actually
\emph{smaller} than on the train split ($0.33$~m S-curve, $0.38$~m
ZUPT on test vs.\ $0.44$~m and $0.56$~m on train respectively), and
the $\le 1.5$~m fraction stays within three percentage points across
the split. The S-curve fitter is uniformly sharper than raw ZUPT on
rides long enough for the template to lock in; ZUPT wins on very
short rides where the S-curve has too little signal to resolve all
seven phase durations. The per-experiment breakdown in
Table~\ref{tab:per-exp-test} shows that two of the three test
experiments achieve $<2$~m MAE on both algorithms — the third, a
single Pixel~10 device at \texttt{beitYitzchaki}, contributes most of
the test-set MAE mass and is the main driver of the gap between the
test median error (very low) and the test mean error (more mixed).

\paragraph{Conformal coverage transfers partially.} Train coverage on
clean data clears the $90\%$ target comfortably ($91.4\%$ ZUPT,
$92.4\%$ S-curve). On the blind test split it drops to $79.6\%$ for
ZUPT and $87.0\%$ for S-curve. The coverage slip is unsurprising in
principle: split-conformal's $1-\alpha$ guarantee is valid under
exchangeability, and a held-out-building split is as clean an
exchangeability violation as one can engineer — the test elevator is
mechanically different from any train elevator, and hand-holding
statistics at \texttt{beitYitzchaki} might differ from the other
locations. The good news is that the residual distribution is not
wildly different from train (median error is actually smaller), so
the failure is concentrated in the \emph{tail} rather than in the
bulk. A post-hoc fix we did \emph{not} take (to preserve the blind-
test protocol) would have been a slightly larger $\alpha$ margin in
the calibration step; the configuration registered in the code uses
$\alpha=0.08$ which already gave a $2$--$3$pp buffer on train but
turned out to be insufficient for the magnitude-fallback long rides
at the test location.

\paragraph{Quality filter and theoretical $\sigma$ analysis.} The
conformal multiplier on the ZUPT side (\texttt{multiplier} field of
\texttt{calibration\_zupt.json}) is approximately $3.7$: after adding
the relative-scale and projection-quality terms the per-segment
$\sigma$ model is about $3\!-\!4\times$ too tight on the worst cases
and conformal absorbs the rest. The S-curve multiplier is $1.80$,
which is consistent with the CRB being a reasonable per-segment
scale once it is safety-factored and combined with the zupt-
disagreement term (see \S\ref{sec:algo-scurve}). The quality filter
rejects $\sim\!13\%$ of clean segments on train and $\sim\!7\%$ on
test; on the small dirty pool (33 segments, all on train) it catches
roughly two-thirds of them, which is the intended asymmetry — the
filter is tuned to be conservative about filtering clean data and
aggressive about flagging dirty data.

\paragraph{Next steps.} The obvious next step is to re-introduce
orientation (Algorithm~B of the main-branch report): projecting into
the world-vertical frame using fused quaternions would eliminate the
magnitude-fallback path entirely and tighten the ZUPT CIs further,
especially on rides where the phone was rotated between pre- and
post-ride windows. That work is deferred to a follow-up since the
\texttt{Predictor} contract currently hands per-segment data without
a quaternion stream. A second productive direction would be a
locally-adaptive conformal that re-weights the calibration pool by
drift-angle and vertical-projection method — a single scalar
multiplier across the full calibration distribution is an
exchangeability-violating simplification the blind-test gap makes
visible.
